% \iffalse
% This document requires lualatex
%%
%% Copyright (C) 2012-2026 Javier Bezos and Johannes L. Braams.
%% Copyright (C) 1989-2012 Johannes L. Braams and
%%           any individual authors listed elsewhere in this file.
%% All rights reserved.
%%
%%
%% This file is part of the Babel system.
%% --------------------------------------
%%
%% It may be distributed and/or modified under the
%% conditions of the LaTeX Project Public License, either version 1.3
%% of this license or (at your option) any later version.
%% The latest version of this license is in
%%   http://www.latex-project.org/lppl.txt
%% and version 1.3 or later is part of all distributions of LaTeX
%% version 2003/12/01 or later.
%%
%% This work has the LPPL maintenance status "maintained".
%%
%% The Current Maintainer of this work is Javier Bezos.
%%
%% The list of derived (unpacked) files belonging to the distribution
%% and covered by LPPL is defined by the unpacking scripts (with
%% extension |.ins|) which are part of the distribution.
%%
% \fi
%
% \iffalse
%<*filedriver>
\ProvidesFile{babel-plain.dtx}%
  [2026/10/18 v26.14.138750
   The multilingual framework for Plain TeX] 
\documentclass{ltxdoc}
\usepackage{babel}
\GetFileInfo{babel-plain.dtx}
\usepackage{fontspec}
\setmainfont[
     Scale=.88,
     BoldFont = NotoSerif-Bold.ttf,
     ItalicFont = NotoSerif-Italic.ttf,
     BoldItalicFont = NotoSerif-BoldItalic.ttf]
  {NotoSerif-Regular.ttf}
\setsansfont[
     Scale=.88,
     BoldFont = NotoSans-Bold.ttf,
     ItalicFont = NotoSans-Italic.ttf,
     BoldItalicFont = NotoSans-BoldItalic.ttf]
  {NotoSans-Regular.ttf}
\setmonofont[Scale=.86, FakeStretch=.97]{DejaVu Sans Mono}
\def\tto{\texttt{‣}}
\def\tlangle{\texttt{⟨}}
\def\trangle{\texttt{⟩}}
\DeclareRobustCommand\meta[1]{\tlangle\mbox{\rmfamily\itshape#1}\trangle}
\raggedright
\setlength{\parindent}{.8em}
\frenchspacing
\addtolength{\oddsidemargin}{-2em}
\addtolength{\textwidth}{25.5pt}
\addtolength{\textheight}{3.5cm}
\addtolength{\topmargin}{-2cm}
\font\manual=logo10 % font used for the METAFONT logo, etc.
\newcommand*\MF{{\manual META}\-{\manual FONT}}
\newcommand*\babel{\textsf{babel}}
\newcommand*\Babel{\textsf{Babel}}
\newcommand*\xetex{\textsf{xetex}}
\newcommand*\pdftex{\textsf{pdftex}}
\newcommand*\luatex{\textsf{luatex}}
\newcommand\largetex{T\kern -.1517em\lower .45ex\hbox {E}\kern -.09emX}
\newcommand*\m[1]{\mbox{{\fontspec{FreeSerif}〈}\normalfont\itshape#1\/}{\fontspec{FreeSerif}〉}}
\newcommand*\langlist{%
  \meta{language}\texttt{,}\meta{language}\texttt{,}...}
\newcommand*\langvar{\m{language}}
\newcommand*\Lopt[1]{\textsf{#1}}
\newcommand*\Lenv[1]{\texttt{#1}}
\newcommand*\menv[1]{\char`\{#1\char`\}}
\newcommand*\Eenv[1]{%
  \quad...\quad
  \texttt{\color{thered}\string\end\menv{#1}}}
\newcommand*\file[1]{\texttt{#1}}
\newcommand*\cls[1]{\texttt{#1}}
\newcommand*\pkg[1]{\texttt{#1}}
\usepackage{multicol}
\usepackage[linkcolor=blue,urlcolor=blue,colorlinks=true,pdfdisplaydoctitle]{hyperref}
\usepackage{tcolorbox}
\tcbuselibrary{breakable}
\usepackage{colortbl}
\hypersetup{%
  pdfsubject={Multilingual documents and localization for pdfLaTeX /
     LuaLaTeX / XeLaTeX},
  pdfauthor={Javier Bezos, Johannes Braams},
  pdftitle={Babel, the multilingual framework to localize LaTeX, LuaLaTeX, and XeLaTeX},
  pdfkeywords={TeX, LaTeX, pdfTeX, LuaTeX, XeTeX, internationalization,
    localization}}
\ExplSyntaxOn
\normalmarginpar
\newcommand\New[1]{%
  \regex_extract_once:nnN{(\d\d?)\.(\d\d?)}{#1}\bbltempa
  \seq_pop_right:NN\bbltempa\bbltempb
  \seq_pop_right:NN\bbltempa\bbltempa
  \leavevmode%\marginpar
  {\scriptsize
  % \enspace
  \colorbox[rgb]{.92, .86, .73}%
    {%New~
     #1%
    \quark_if_no_value:NTF\bbltempb{}%
      {~
       \ifnum\bbltempa>2 \ifnum\bbltempb>34
         \href{https://latex3.github.io/babel/news/%
               whats-new-in-babel-\bbltempa.\bbltempb.html}%
              {\scriptsize ※}%
       \fi\fi
       \ifnum\bbltempa>23
         \href{https://latex3.github.io/babel/news/%
               whats-new-in-babel-\bbltempa.\bbltempb.html}%
              {\scriptsize ※}%
       \fi}}}\enspace\ignorespaces}
\ExplSyntaxOff
\definecolor{thered}{rgb}{0.65,0.04,0.07}
\definecolor{thegrey}{gray}{0.8}
\definecolor{thegreyed}{gray}{0.6}
\definecolor{shadecolor}{rgb}{1,1,0.97}
\definecolor{messages}{rgb}{.66,.13,.27}
\makeatletter
\def\@seccntformat#1{%
  \llap{\makebox[1.7cm][l]{\csname the#1\endcsname.}}%
  \hspace{-.7cm}}
\newcommand\displayed[1]{%
  \setlength{\leftmargini}{1.5em}%
  \list{}{}\item[]%
  \hskip-\leftmargin{\sffamily\bfseries\color{messages}\MakeUppercase{#1\ }}}
\makeatother
\newenvironment{warning}{\displayed{Warning}}{\endlist}
\newenvironment{note}{\displayed{Note}}{\endlist}
\newenvironment{example}{\displayed{Example}}{\endlist}
\newenvironment{troubleshooting}{\displayed{Troubleshooting}}{\endlist}
\newcommand\setengine{\def\engine}
\let\engine\relax
\begingroup
  \catcode`\<=13
  \catcode`\>=13
  \catcode`\|=13
  \AtBeginDocument{%
    \gdef|{\verb|\def<##1>{{\fontspec{FreeSerif}〈}{\rmfamily\itshape##1}{\fontspec{FreeSerif}〉}}}}
\endgroup
\AddToHook{env/verbatim/before}{%
  \begin{tcolorbox}[breakable, 
    sharp corners, colframe=black!50,
    boxrule=.4pt, bottomrule at break = 0pt, toprule at break = 0pt, 
    left=1.5mm, right=1.5mm, top=1.5mm, bottom=1.5mm,
    colback=shadecolor]%
  \small
}
\AddToHook{env/verbatim/after}{%
  \end{tcolorbox}%
}
\catcode`×=\active
\def×{\bgroup\let×\egroup\leavevmode\color{thered}}
\def\MacroFont{\fontencoding \encodingdefault \fontfamily\ttdefault
  \fontseries\mddefault \fontshape\updefault \small \catcode`\_=\active}
\definecolor{shadecolor}{rgb}{0.97,0.97,0.94}
%
\makeatletter   
\def\Describe#1#2{%
  \vskip 1.7ex\@plus .7ex \@minus .2ex
  \noindent\hspace*{-1.7cm}%
  {\color{thegreyed}\normalsize\ttfamily
   \underline{\strut{\color{thered}\string#1}\color{black}#2}}%
  \par\nobreak\vskip 1.1ex \@plus .2ex }
\def\DescribeOther{\nobreak\vskip-1.1\baselineskip\Describe}
\AtBeginDocument{\def\DescribeMacro#1{{\ttfamily\color{thered}\string#1}}}
\def\macro#1{%
  \if@noskipsec\leavevmode\vskip-1.2\baselineskip\vskip0pt\fi
  \@startsection{paragraph}{4}{\z@}%
    {3.25ex \@plus1ex \@minus.2ex}%
    {-1em}%
    {\sffamily\bfseries\color{thered}}*{\hspace*{-1em}\string#1}}
\let\environment\macro
\def\optstar{%
  \kern1pt\relax
  \colorbox{thegrey}{\ttfamily\hskip-.2em*\hskip-.2em}%
  \kern1pt}
\def\trouble#1{\addcontentsline{tsh}{trouble}{#1}}
\def\toclevel@trouble{0}
\def\listoftroubles{\section*{Troubleshooting}\@starttoc{tsh}}
\let\l@trouble\l@figure
\def\@tocrmarg{3.5em plus 1000em}
\let\saved@check@percent\check@percent
\let\check@percent\relax
% Changes to doc.sty (add <<...>> syntax)
\def\Module#1{%
  \mod@math@codes\tlangle{\color{thered}\textsf{#1}}\trangle}
\begingroup 
  \catcode`\<\active
  \catcode`\>\active
  \gdef\check@plus@etc{%
    \let\bbl@next\pm@module
    \ifx*\next
      \let\bbl@next\star@module
    \else\ifx/\next
      \let\bbl@next\slash@module
    \else\ifx<\next
      \let\bbl@next\var@module
    \fi\fi\fi
    \bbl@next}
  \gdef\var@module#1#2#3>>{%
     \tlangle\!\pm@module#2#3>\!\textup{\trangle}
     \ifx*#2\ \tto\fi}
\endgroup
\renewcommand\partname{Part}
\def\@pnumwidth{3em}
\makeatother
\begin{document}
\title{Babel, the multilingual framework to localize
   LaTeX, LuaLaTeX, and XeLaTeX}
\author{Johannes Braams\\
        Javier Bezos}
\date{Typeset \today}
\begin{titlepage}
\begin{picture}[artifact](0,0)
  \sffamily
  \put(-5.7cm,-4.5cm){\color[rgb]{.86,.73,.67}\rule{8cm}{8cm}}
  \put( 2.3cm,-4.5cm){\color[rgb]{.95,.95,.95}\rule{15cm}{8cm}}
  \put(-5.5cm,-27.5cm){\color[rgb]{.95,.95,.95}\rule{8cm}{23cm}}
  \put(2.3cm,-27.5cm){\color[rgb]{.92, .86, .73}\rule{15cm}{23cm}}
  \put(3cm,-2.6cm){\color[rgb]{.55,.4,.35}\fontsize{72}{72}\selectfont Babel}
  \put(-4.1cm,-10.8cm){\begin{minipage}{5cm}
    \Huge
    \strut

    \Large
    \vspace{1.2cm}
    Version \csname @gobble\expandafter\endcsname\fileversion\newline
    \filedate

    \LARGE
    \vspace{1.2cm}
    Javier Bezos\\[-.5ex]
    {\large Current maintainer}

    \vspace{.3cm}
    Johannes L. Braams\\[-.5ex]
    {\large Original author}
    \end{minipage}}
  \put( 3cm, -12.8cm){\begin{minipage}{10cm}
    \fontsize{35}{45}\selectfont
    \setlength\parskip{3mm}\raggedright
    Localization and internationalization\\[1cm]
    Plain \largetex
    \end{minipage}}
\end{picture}
\end{titlepage}

\begin{note}
  This file is devoted to \babel{} for Plain \TeX. For many years, most
  of the code was shared between \LaTeX{} and Plain, but recent
  advances in \LaTeX{} have made joint maintenance increasingly
  complex. Consequently, the code for Plain will be split from main
  \babel{} for \LaTeX{} and essentially frozen. This does not mean it
  is unmaintained, and the door remains open to future improvements.
  
  Description of many macros are the same as those in
  \textsf{babel-code.pdf}, and here only some hints are provided.
  Please, refer to that document for details.
  
  This file is loaded not only by Plain. \LaTeX{} loads it with the
  package option \texttt{base} (and also \texttt{plain.def}) with
  |\bbl@onlyswitch| $=$ |\@empty|. See sec.~\ref{endswitch}.
  
  It needs to be cleaned up to remove unnecessary code in Plain. Note
  also the engine specific files are still loaded.
\end{note}

\DocInput{babel-plain.dtx}

\end{document}
%</filedriver>
% \fi
%
% \addtolength{\oddsidemargin}{-5pt}
% \addtolength{\textwidth}{50pt}
% \reversemarginpar
%
% \section{A few preliminary definitions}
%
% Just a block. To be revised. 
%
%    \begin{macrocode}
%<<version=26.14.138750>>
%<<date=2026/10/18>>
%    \end{macrocode}
%
% Some files identify themselves with a \LaTeX{} macro. The following
% code is placed before them to define (and then undefine) if not in
% \LaTeX.
%
% \macro{\language}
%
% Just for compatibility, for not to touch |hyphen.cfg|.
%
%    \begin{macrocode}
%<<*Define Plain core switching macros>>
\ifx\language\@undefined
  \csname newcount\endcsname\language
\fi
%<</Define Plain core switching macros>>
%    \end{macrocode}
%
% \macro{\last@language}
%
% Another counter is used to keep track of the allocated languages.
% \TeX{} and \LaTeX{} reserves for this purpose the count 19.
%
% \macro{\addlanguage}
%
% This macro was introduced for \TeX{} $<$ 2. Preserved for
% compatibility.
%
%    \begin{macrocode}
%<<*Define Plain core switching macros>>
\countdef\last@language=19
\def\addlanguage{\csname newlanguage\endcsname}
%<</Define Plain core switching macros>>
%    \end{macrocode}
%
% First, exit immediately if previouly loaded.
%
%    \begin{macrocode}
%<*core>
\ifx\ldf@quit\@undefined\else
\endinput\fi % Same line!
\ifx\ProvidesFile\@undefined
  \def\ProvidesFile#1[#2 #3 #4]{%
    \wlog{File: #1 #4 #3 <#2>}%
    \let\ProvidesFile\@undefined}
\fi
\ProvidesFile{babel.def}[<@date@> v<@version@> Babel common definitions]
%    \end{macrocode}
%
% \section{Emulating some \LaTeX{} features}
%
% The file |babel.def| expects some definitions made in the \LaTeXe{}
% style file. So, in Plain{} we must provide at least
% some predefined values as well some tools to set them (even if not
% all options are available). There are no package options, and
% therefore and alternative mechanism is provided. For the moment,
% only |\babeloptionstrings| and |\babeloptionmath| are provided,
% which can be defined before loading \babel. |\BabelModifiers| can be
% set too (but not sure it works).
%  
%    \begin{macrocode}
\ifx\AtBeginDocument\@undefined
\def\@empty{}
\def\loadlocalcfg#1{%
  \openin0#1.cfg
  \ifeof0
    \closein0
  \else
    \closein0
    {\immediate\write16{*************************************}%
     \immediate\write16{* Local config file #1.cfg used}%
     \immediate\write16{*}%
     }
    \input #1.cfg\relax
  \fi
  \@endofldf}
%    \end{macrocode}
%
% \subsection{General tools}
%
% A number of \LaTeX\ macro's that are needed later on.
%
%    \begin{macrocode}
\long\def\@firstofone#1{#1}
\long\def\@firstoftwo#1#2{#1}
\long\def\@secondoftwo#1#2{#2}
\def\@nnil{\@nil}
\def\@gobbletwo#1#2{}
\def\@ifstar#1{\@ifnextchar *{\@firstoftwo{#1}}}
\def\@star@or@long#1{%
  \@ifstar
  {\let\l@ngrel@x\relax#1}%
  {\let\l@ngrel@x\long#1}}
\let\l@ngrel@x\relax
\def\@car#1#2\@nil{#1}
\def\@cdr#1#2\@nil{#2}
\let\@typeset@protect\relax
\let\protected@edef\edef
\long\def\@gobble#1{}
\edef\@backslashchar{\expandafter\@gobble\string\\}
\def\strip@prefix#1>{}
\def\g@addto@macro#1#2{{%
    \toks@\expandafter{#1#2}%
    \xdef#1{\the\toks@}}}
\def\@namedef#1{\expandafter\def\csname #1\endcsname}
\def\@nameuse#1{\csname #1\endcsname}
\def\@ifundefined#1{%
  \expandafter\ifx\csname#1\endcsname\relax
    \expandafter\@firstoftwo
  \else
    \expandafter\@secondoftwo
  \fi}
\def\@expandtwoargs#1#2#3{%
  \edef\reserved@a{\noexpand#1{#2}{#3}}\reserved@a}
\def\zap@space#1 #2{%
  #1%
  \ifx#2\@empty\else\expandafter\zap@space\fi
  #2}
\let\bbl@trace\@gobble
\def\bbl@error#1{% Implicit #2#3#4
  \begingroup
    \catcode`\\=0   \catcode`\==12 \catcode`\`=12
    \catcode`\^^M=5 \catcode`\%=14
    \input errbabel.def
  \endgroup
  \bbl@error{#1}}
\def\bbl@warning#1{%
  \begingroup
    \newlinechar=`\^^J
    \def\\{^^J(babel) }%
    \message{\\#1}%
  \endgroup}
\let\bbl@infowarn\bbl@warning
\def\bbl@info#1{%
  \begingroup
    \newlinechar=`\^^J
    \def\\{^^J}%
    \wlog{#1}%
  \endgroup}
%    \end{macrocode}
%
%    \LaTeXe\ has the command |\@onlypreamble| which adds commands to
%    a list of commands that are no longer needed after
%    |\begin{document}|.
%
%    \begin{macrocode}
\ifx\@preamblecmds\@undefined
  \def\@preamblecmds{}
\fi
\def\@onlypreamble#1{%
  \expandafter\gdef\expandafter\@preamblecmds\expandafter{%
    \@preamblecmds\do#1}}
\@onlypreamble\@onlypreamble
%    \end{macrocode}
%
% Mimic \LaTeX's |\AtBeginDocument|; for this to work the user needs to
% add |\begindocument| to his file.
%
%    \begin{macrocode}
\def\begindocument{%
  \@begindocumenthook
  \global\let\@begindocumenthook\@undefined
  \def\do##1{\global\let##1\@undefined}%
  \@preamblecmds
  \global\let\do\noexpand}
%    \end{macrocode}
%
%    \begin{macrocode}
\ifx\@begindocumenthook\@undefined
  \def\@begindocumenthook{}
\fi
\@onlypreamble\@begindocumenthook
\def\AtBeginDocument{\g@addto@macro\@begindocumenthook}
%    \end{macrocode}
%
% We also have to mimic \LaTeX's |\AtEndOfPackage|. Our replacement
% macro is much simpler; it stores its argument in |\@endofldf|.
%
%    \begin{macrocode}
\def\AtEndOfPackage#1{\g@addto@macro\@endofldf{#1}}
\@onlypreamble\AtEndOfPackage
\def\@endofldf{}
\@onlypreamble\@endofldf
\let\bbl@afterlang\@empty
\chardef\bbl@opt@hyphenmap\z@
%    \end{macrocode}
%
% \LaTeX\ needs to be able to switch off writing to its auxiliary
% files; plain doesn't have them by default. There is a trick to hide
% some conditional commands from the outer |\ifx|. The same trick is
% applied below.
%
%    \begin{macrocode}
\catcode`\&=\z@
\ifx&if@filesw\@undefined
  \expandafter\let\csname if@filesw\expandafter\endcsname
    \csname iffalse\endcsname
\fi
\catcode`\&=4
%    \end{macrocode}
%
% Mimic \LaTeX's commands to define control sequences.
%
%    \begin{macrocode}
\def\newcommand{\@star@or@long\new@command}
\def\new@command#1{%
  \@testopt{\@newcommand#1}0}
\def\@newcommand#1[#2]{%
  \@ifnextchar [{\@xargdef#1[#2]}%
                {\@argdef#1[#2]}}
\long\def\@argdef#1[#2]#3{%
  \@yargdef#1\@ne{#2}{#3}}
\long\def\@xargdef#1[#2][#3]#4{%
  \expandafter\def\expandafter#1\expandafter{%
    \expandafter\@protected@testopt\expandafter #1%
    \csname\string#1\expandafter\endcsname{#3}}%
  \expandafter\@yargdef \csname\string#1\endcsname
  \tw@{#2}{#4}}
\long\def\@yargdef#1#2#3{%
  \@tempcnta#3\relax
  \advance \@tempcnta \@ne
  \let\@hash@\relax
  \edef\reserved@a{\ifx#2\tw@ [\@hash@1]\fi}%
  \@tempcntb #2%
  \@whilenum\@tempcntb <\@tempcnta
  \do{%
    \edef\reserved@a{\reserved@a\@hash@\the\@tempcntb}%
    \advance\@tempcntb \@ne}%
  \let\@hash@##%
  \l@ngrel@x\expandafter\def\expandafter#1\reserved@a}
\def\providecommand{\@star@or@long\provide@command}
\def\provide@command#1{%
  \begingroup
    \escapechar\m@ne\xdef\@gtempa{{\string#1}}%
  \endgroup
  \expandafter\@ifundefined\@gtempa
    {\def\reserved@a{\new@command#1}}%
    {\let\reserved@a\relax
     \def\reserved@a{\new@command\reserved@a}}%
   \reserved@a}%
%    \end{macrocode}
%
%    \begin{macrocode}
\def\DeclareRobustCommand{\@star@or@long\declare@robustcommand}
\def\declare@robustcommand#1{%
   \edef\reserved@a{\string#1}%
   \def\reserved@b{#1}%
   \edef\reserved@b{\expandafter\strip@prefix\meaning\reserved@b}%
   \edef#1{%
      \ifx\reserved@a\reserved@b
         \noexpand\x@protect
         \noexpand#1%
      \fi
      \noexpand\protect
      \expandafter\noexpand\csname
         \expandafter\@gobble\string#1 \endcsname
   }%
   \expandafter\new@command\csname
      \expandafter\@gobble\string#1 \endcsname
}
\def\x@protect#1{%
   \ifx\protect\@typeset@protect\else
      \@x@protect#1%
   \fi
}
\catcode`\&=\z@  % Trick to hide conditionals
  \def\@x@protect#1&fi#2#3{&fi\protect#1}
%    \end{macrocode}
%
% The following little macro |\in@| is taken from \file{latex.ltx}; it
% checks whether its first argument is part of its second argument. It
% uses the boolean |\in@|; allocating a new boolean inside
% conditionally executed code is not possible, hence the construct with
% the temporary definition of |\bbl@tempa|.
%
%    \begin{macrocode}
  \def\bbl@tempa{\csname newif\endcsname&ifin@}
\catcode`\&=4
\ifx\in@\@undefined
  \def\in@#1#2{%
    \def\in@@##1#1##2##3\in@@{%
      \ifx\in@##2\in@false\else\in@true\fi}%
    \in@@#2#1\in@\in@@}
\else
  \let\bbl@tempa\@empty
\fi
\bbl@tempa
%    \end{macrocode}
%
% Used in some \textsf{ldf} files. Assume true.
%
%    \begin{macrocode}
\def\@ifpackagewith#1#2#3#4{#3}
%    \end{macrocode}
%
% The \LaTeX\ macro |\@ifl@aded| checks whether a file was loaded.
% No-op in Plain.
%
%    \begin{macrocode}
\def\@ifl@aded#1#2#3#4{}
%    \end{macrocode}
%
% For the following code we need to make sure that the commands
% |\newcommand| and |\providecommand| exist with some sensible
% definition. They are not fully equivalent to their \LaTeXe\ versions;
% just enough to make things work in plain~\TeX environments.
%
%    \begin{macrocode}
\ifx\@tempcnta\@undefined
  \csname newcount\endcsname\@tempcnta\relax
\fi
\ifx\@tempcntb\@undefined
  \csname newcount\endcsname\@tempcntb\relax
\fi
%    \end{macrocode}
%
% To prevent wasting two counters in \LaTeX\ (because counters with the
% same name are allocated later by it) we reset the counter that holds
% the next free counter (|\count10|).
%
%    \begin{macrocode}
\ifx\bye\@undefined
  \advance\count10 by -2\relax
\fi
\ifx\@ifnextchar\@undefined
  \def\@ifnextchar#1#2#3{%
    \let\reserved@d=#1%
    \def\reserved@a{#2}\def\reserved@b{#3}%
    \futurelet\@let@token\@ifnch}
  \def\@ifnch{%
    \ifx\@let@token\@sptoken
      \let\reserved@c\@xifnch
    \else
      \ifx\@let@token\reserved@d
        \let\reserved@c\reserved@a
      \else
        \let\reserved@c\reserved@b
      \fi
    \fi
    \reserved@c}
  \def\:{\let\@sptoken= } \:  % this makes \@sptoken a space token
  \def\:{\@xifnch} \expandafter\def\: {\futurelet\@let@token\@ifnch}
\fi
\def\@testopt#1#2{%
  \@ifnextchar[{#1}{#1[#2]}}
\def\@protected@testopt#1{%
  \ifx\protect\@typeset@protect
    \expandafter\@testopt
  \else
    \@x@protect#1%
  \fi}
\long\def\@whilenum#1\do #2{\ifnum #1\relax #2\relax\@iwhilenum{#1\relax
     #2\relax}\fi}
\long\def\@iwhilenum#1{\ifnum #1\expandafter\@iwhilenum
         \else\expandafter\@gobble\fi{#1}}
%    \end{macrocode}
%
% \subsection{Encoding related macros}
%
% Code from \file{ltoutenc.dtx}, adapted for use in the plain \TeX\
% environment.
%
%    \begin{macrocode}
\def\DeclareTextCommand{%
   \@dec@text@cmd\providecommand
}
\def\ProvideTextCommand{%
   \@dec@text@cmd\providecommand
}
\def\DeclareTextSymbol#1#2#3{%
   \@dec@text@cmd\chardef#1{#2}#3\relax
}
\def\@dec@text@cmd#1#2#3{%
   \expandafter\def\expandafter#2%
      \expandafter{%
         \csname#3-cmd\expandafter\endcsname
         \expandafter#2%
         \csname#3\string#2\endcsname
      }%
%   \let\@ifdefinable\@rc@ifdefinable
   \expandafter#1\csname#3\string#2\endcsname
}
\def\@current@cmd#1{%
  \ifx\protect\@typeset@protect\else
      \noexpand#1\expandafter\@gobble
  \fi
}
\def\@changed@cmd#1#2{%
   \ifx\protect\@typeset@protect
      \expandafter\ifx\csname\cf@encoding\string#1\endcsname\relax
         \expandafter\ifx\csname ?\string#1\endcsname\relax
            \expandafter\def\csname ?\string#1\endcsname{%
               \@changed@x@err{#1}%
            }%
         \fi
         \global\expandafter\let
           \csname\cf@encoding \string#1\expandafter\endcsname
           \csname ?\string#1\endcsname
      \fi
      \csname\cf@encoding\string#1%
        \expandafter\endcsname
   \else
      \noexpand#1%
   \fi
}
\def\@changed@x@err#1{%
    \errhelp{Your command will be ignored, type <return> to proceed}%
    \errmessage{Command \protect#1 undefined in encoding \cf@encoding}}
\def\DeclareTextCommandDefault#1{%
   \DeclareTextCommand#1?%
}
\def\ProvideTextCommandDefault#1{%
   \ProvideTextCommand#1?%
}
\expandafter\let\csname OT1-cmd\endcsname\@current@cmd
\expandafter\let\csname?-cmd\endcsname\@changed@cmd
\def\DeclareTextAccent#1#2#3{%
  \DeclareTextCommand#1{#2}[1]{\accent#3 ##1}
}
\def\DeclareTextCompositeCommand#1#2#3#4{%
   \expandafter\let\expandafter\reserved@a\csname#2\string#1\endcsname
   \edef\reserved@b{\string##1}%
   \edef\reserved@c{%
     \expandafter\@strip@args\meaning\reserved@a:-\@strip@args}%
   \ifx\reserved@b\reserved@c
      \expandafter\expandafter\expandafter\ifx
         \expandafter\@car\reserved@a\relax\relax\@nil
         \@text@composite
      \else
         \edef\reserved@b##1{%
            \def\expandafter\noexpand
               \csname#2\string#1\endcsname####1{%
               \noexpand\@text@composite
                  \expandafter\noexpand\csname#2\string#1\endcsname
                  ####1\noexpand\@empty\noexpand\@text@composite
                  {##1}%
            }%
         }%
         \expandafter\reserved@b\expandafter{\reserved@a{##1}}%
      \fi
      \expandafter\def\csname\expandafter\string\csname
         #2\endcsname\string#1-\string#3\endcsname{#4}
   \else
     \errhelp{Your command will be ignored, type <return> to proceed}%
     \errmessage{\string\DeclareTextCompositeCommand\space used on
         inappropriate command \protect#1}
   \fi
}
\def\@text@composite#1#2#3\@text@composite{%
   \expandafter\@text@composite@x
      \csname\string#1-\string#2\endcsname
}
\def\@text@composite@x#1#2{%
   \ifx#1\relax
      #2%
   \else
      #1%
   \fi
}
%
\def\@strip@args#1:#2-#3\@strip@args{#2}
\def\DeclareTextComposite#1#2#3#4{%
   \def\reserved@a{\DeclareTextCompositeCommand#1{#2}{#3}}%
   \bgroup
      \lccode`\@=#4%
      \lowercase{%
   \egroup
      \reserved@a @%
   }%
}
%
\def\UseTextSymbol#1#2{#2}
\def\UseTextAccent#1#2#3{}
\def\@use@text@encoding#1{}
\def\DeclareTextSymbolDefault#1#2{%
   \DeclareTextCommandDefault#1{\UseTextSymbol{#2}#1}%
}
\def\DeclareTextAccentDefault#1#2{%
   \DeclareTextCommandDefault#1{\UseTextAccent{#2}#1}%
}
\def\cf@encoding{OT1}
%    \end{macrocode}
%
% Currently we only use the \LaTeXe\ method for accents for those that
% are known to be made active in \emph{some} language definition file.
%
%    \begin{macrocode}
\DeclareTextAccent{\"}{OT1}{127}
\DeclareTextAccent{\'}{OT1}{19}
\DeclareTextAccent{\^}{OT1}{94}
\DeclareTextAccent{\`}{OT1}{18}
\DeclareTextAccent{\~}{OT1}{126}
%    \end{macrocode}
%
% The following control sequences are used in \file{babel.def} but are
% not defined for \textsc{plain} \TeX.
%
%    \begin{macrocode}
\DeclareTextSymbol{\textquotedblleft}{OT1}{92}
\DeclareTextSymbol{\textquotedblright}{OT1}{`\"}
\DeclareTextSymbol{\textquoteleft}{OT1}{`\`}
\DeclareTextSymbol{\textquoteright}{OT1}{`\'}
\DeclareTextSymbol{\i}{OT1}{16}
\DeclareTextSymbol{\ss}{OT1}{25}
%    \end{macrocode}
%
% For a couple of languages we need the \LaTeX-control sequence
% |\scriptsize| to be available. Because plain \TeX\ doesn't have such
% a sophisticated font mechanism as \LaTeX\ has, we just |\let| it to
% |\sevenrm|.
%
%    \begin{macrocode}
\ifx\scriptsize\@undefined
  \let\scriptsize\sevenrm
\fi
%    \end{macrocode}
%
% And a few more “dummy” definitions.
%
%    \begin{macrocode}
\def\languagename{english}%
\def\bbl@ifshorthand#1#2#3{#2}%
\let\bbl@language@opts\@empty
\let\bbl@provide@locale\relax
\ifx\babeloptionstrings\@undefined
  \let\bbl@opt@strings\@nnil
\else
  \let\bbl@opt@strings\babeloptionstrings
\fi
\def\BabelStringsDefault{generic}
\def\bbl@tempa{normal}
\ifx\babeloptionmath\bbl@tempa
  \def\bbl@mathnormal{\noexpand\textormath}
\fi
\def\AfterBabelLanguage#1#2{}
\ifx\BabelModifiers\@undefined\let\BabelModifiers\relax\fi
\let\bbl@afterlang\relax
\ifx\@uclclist\@undefined\let\@uclclist\@empty\fi
\chardef\bbl@bidimode\z@
\fi
%    \end{macrocode}
%
% \section{Basic macros}
%
% The emulation is finished and now comes the basic macros for Plain +
% \texttt{base}. See \textsf{babel-code.pdf}.
%
%    \begin{macrocode}
\bbl@trace{Basic macros}
\def\bbl@stripslash{\expandafter\@gobble\string}
\def\bbl@add#1#2{%
  \bbl@ifunset{\bbl@stripslash#1}%
    {\def#1{#2}}%
    {\expandafter\def\expandafter#1\expandafter{#1#2}}}
\def\bbl@xin@{\@expandtwoargs\in@}
\def\bbl@carg#1#2{\expandafter#1\csname#2\endcsname}%
\def\bbl@ncarg#1#2#3{\expandafter#1\expandafter#2\csname#3\endcsname}%
\def\bbl@ccarg#1#2#3{%
  \expandafter#1\csname#2\expandafter\endcsname\csname#3\endcsname}%
\def\bbl@csarg#1#2{\expandafter#1\csname bbl@#2\endcsname}%
\def\bbl@cs#1{\csname bbl@#1\endcsname}
\def\bbl@cl#1{\csname bbl@#1@\languagename\endcsname}
\def\bbl@loop#1#2#3{\bbl@@loop#1{#3}#2,\@nnil,}
\def\bbl@loopx#1#2{\expandafter\bbl@loop\expandafter#1\expandafter{#2}}
\def\bbl@@loop#1#2#3,{%
  \ifx\@nnil#3\relax\else
    \def#1{#3}#2\bbl@afterfi\bbl@@loop#1{#2}%
  \fi}
\def\bbl@for#1#2#3{\bbl@loopx#1{#2}{\ifx#1\@empty\else#3\fi}}
%    \end{macrocode}
%
% \macro{\bbl@add@list}
% This internal macro adds its second argument to a comma separated
% list in its first argument.
%
%    \begin{macrocode}
\def\bbl@add@list#1#2{%
  \edef#1{%
    \bbl@ifunset{\bbl@stripslash#1}%
      {}%
      {\ifx#1\@empty\else#1,\fi}%
    #2}}
%    \end{macrocode}
%
% \macro{\bbl@afterelse}
% \macro{\bbl@afterfi}
% Because the code that is used in the handling of active characters
% may need to look ahead, we take extra care to `throw' it over the
% |\else| and |\fi| parts of an |\if|-statement.
%
%    \begin{macrocode}
\long\def\bbl@afterelse#1\else#2\fi{\fi#1}
\long\def\bbl@afterfi#1\fi{\fi#1}
%    \end{macrocode}
%
% \macro{\bbl@exp}
%
% Now, just syntactical sugar.
%
%    \begin{macrocode}
\def\bbl@exp#1{%
  \begingroup
    \let\\\noexpand
    \let\<\bbl@exp@en
    \let\[\bbl@exp@ue
    \edef\bbl@exp@aux{\endgroup#1}%
  \bbl@exp@aux}
\def\bbl@exp@en#1>{\expandafter\noexpand\csname#1\endcsname}%
\def\bbl@exp@ue#1]{%
  \unexpanded\expandafter\expandafter\expandafter{\csname#1\endcsname}}%
%    \end{macrocode}
%
% \macro{\bbl@trim}
% 
% Define two macros: |\bbl@trim| and |\bbl@trim@def|.
%
%    \begin{macrocode}
\def\bbl@tempa#1{%
  \long\def\bbl@trim##1##2{%
    \futurelet\bbl@trim@a\bbl@trim@c##2\@nil\@nil#1\@nil\relax{##1}}%
  \def\bbl@trim@c{%
    \ifx\bbl@trim@a\@sptoken
      \expandafter\bbl@trim@b
    \else
      \expandafter\bbl@trim@b\expandafter#1%
    \fi}%
  \long\def\bbl@trim@b#1##1 \@nil{\bbl@trim@i##1}}
\bbl@tempa{ }
\long\def\bbl@trim@i#1\@nil#2\relax#3{#3{#1}}
\long\def\bbl@trim@def#1{\bbl@trim{\def#1}}
%    \end{macrocode}
%
% \macro{\bbl@ifunset}
% To check if a macro is defined, we create a new macro, which does the
% same as |\@ifundefined|. However, in an $\varepsilon$-tex engine, it is
% based on |\ifcsname|, which is more efficient, and does not waste
% memory. Defined inside a group, to avoid \cs{ifcsname} being
% implicitly set to |\relax| by the |\csname| test.
%
%    \begin{macrocode}
\begingroup
  \gdef\bbl@ifunset#1{%
    \expandafter\ifx\csname#1\endcsname\relax
      \expandafter\@firstoftwo
    \else
      \expandafter\@secondoftwo
    \fi}
  \bbl@ifunset{ifcsname}%
    {}%
    {\gdef\bbl@ifunset#1{%
       \ifcsname#1\endcsname
         \expandafter\ifx\csname#1\endcsname\relax
           \bbl@afterelse\expandafter\@firstoftwo
         \else
           \bbl@afterfi\expandafter\@secondoftwo
         \fi
       \else
         \expandafter\@firstoftwo
       \fi}}
\endgroup
%    \end{macrocode}
%
% \macro{\bbl@ifblank}
% A tool from \textsf{url}, by Donald Arseneau, which tests if a string
% is empty or space. 
%
%    \begin{macrocode}
\def\bbl@ifblank#1{%
  \bbl@ifblank@i#1\@nil\@nil\@secondoftwo\@firstoftwo\@nil}
\long\def\bbl@ifblank@i#1#2\@nil#3#4#5\@nil{#4}
\def\bbl@ifset#1#2#3{%
  \bbl@ifunset{#1}{#3}{\bbl@exp{\\\bbl@ifblank{\@nameuse{#1}}}{#3}{#2}}}
%    \end{macrocode}
%
% For each element in the comma separated <key>|=|<value> list,
% execute <code> with |#1| and |#2| as the key and the value of
% current item (trimmed).
%
%    \begin{macrocode}
\def\bbl@forkv#1#2{%
  \def\bbl@kvcmd##1##2##3{#2}%
  \bbl@kvnext#1,\@nil,}
\def\bbl@kvnext#1,{%
  \ifx\@nil#1\relax\else
    \bbl@ifblank{#1}{}{\bbl@forkv@eq#1=\@empty=\@nil{#1}}%
    \expandafter\bbl@kvnext
  \fi}
\def\bbl@forkv@eq#1=#2=#3\@nil#4{%
  \bbl@trim@def\bbl@forkv@a{#1}%
  \bbl@trim{\expandafter\bbl@kvcmd\expandafter{\bbl@forkv@a}}{#2}{#4}}
%    \end{macrocode}
%
% A \textit{for} loop. Each item (trimmed) is |#1|.  It cannot be
% nested (it’s doable, but we don't need it).
%
%    \begin{macrocode}
\def\bbl@vforeach#1#2{%
  \def\bbl@forcmd##1{#2}%
  \bbl@fornext#1,\@nil,}
\def\bbl@fornext#1,{%
  \ifx\@nil#1\relax\else
    \bbl@ifblank{#1}{}{\bbl@trim\bbl@forcmd{#1}}%
    \expandafter\bbl@fornext
  \fi}
\def\bbl@foreach#1{\expandafter\bbl@vforeach\expandafter{#1}}
%    \end{macrocode}
%
% Some code should be executed once. The first argument is a flag.
%
%    \begin{macrocode}
\global\let\bbl@done\@empty
\def\bbl@once#1#2{%
  \bbl@xin@{,#1,}{,\bbl@done,}%
  \ifin@\else
    #2%
    \xdef\bbl@done{\bbl@done,#1,}%
  \fi}
%    \end{macrode}
%
% \macro{\bbl@replace}
%
% Returns implicitly |\toks@| with the modified string. 
%
%    \begin{macrocode}
\def\bbl@replace#1#2#3{% in #1 -> repl #2 by #3
  \toks@{}%
  \def\bbl@replace@aux##1#2##2#2{%
    \ifx\bbl@nil##2%
      \toks@\expandafter{\the\toks@##1}%
    \else
      \toks@\expandafter{\the\toks@##1#3}%
      \bbl@afterfi
      \bbl@replace@aux##2#2%
    \fi}%
  \expandafter\bbl@replace@aux#1#2\bbl@nil#2%
  \edef#1{\the\toks@}}
%    \end{macrocode}
%
% An extension to the previous macro:
%
%    \begin{macrocode}
\ifx\detokenize\@undefined\else % Unused macros if old Plain TeX
  \bbl@exp{\def\\\bbl@parsedef##1\detokenize{macro:}}#2->#3\relax{%
    \def\bbl@tempa{#1}%
    \def\bbl@tempb{#2}%
    \def\bbl@tempe{#3}}
  \def\bbl@sreplace#1#2#3{%
    \begingroup
      \expandafter\bbl@parsedef\meaning#1\relax
      \def\bbl@tempc{#2}%
      \edef\bbl@tempc{\expandafter\strip@prefix\meaning\bbl@tempc}%
      \def\bbl@tempd{#3}%
      \edef\bbl@tempd{\expandafter\strip@prefix\meaning\bbl@tempd}%
      \bbl@xin@{\bbl@tempc}{\bbl@tempe}% If not in macro, do nothing
      \ifin@
        \bbl@exp{\\\bbl@replace\\\bbl@tempe{\bbl@tempc}{\bbl@tempd}}%
        \def\bbl@tempc{%     Expanded an executed below as 'uplevel'
           \\\makeatletter % "internal" macros with @ are assumed
           \\\scantokens{%
             \bbl@tempa\\\@namedef{\bbl@stripslash#1}\bbl@tempb{\bbl@tempe}%
             \noexpand\noexpand}%
           \catcode64=\the\catcode64\relax}%  Restore @
      \else
        \let\bbl@tempc\@empty  % Not \relax
      \fi
      \bbl@exp{%      For the 'uplevel' assignments
    \endgroup
      \bbl@tempc}}  % empty or expand to set #1 with changes 
\fi
%    \end{macrocode}
%
% Two further tools. 
%
%    \begin{macrocode}
\def\bbl@ifsamestring#1#2{%
  \begingroup
    \protected@edef\bbl@tempb{#1}%
    \edef\bbl@tempb{\expandafter\strip@prefix\meaning\bbl@tempb}%
    \protected@edef\bbl@tempc{#2}%
    \edef\bbl@tempc{\expandafter\strip@prefix\meaning\bbl@tempc}%
    \ifx\bbl@tempb\bbl@tempc
      \aftergroup\@firstoftwo
    \else
      \aftergroup\@secondoftwo
    \fi
  \endgroup}
\chardef\bbl@engine=%
  \ifx\directlua\@undefined
    \ifx\XeTeXinputencoding\@undefined
      \z@
    \else
      \tw@
    \fi
  \else
    \@ne
  \fi
%    \end{macrocode}
%
% A somewhat hackish tool (hence its name) to avoid spurious spaces in
% some contexts.
%
%    \begin{macrocode}
\def\bbl@bsphack{%
  \ifhmode
    \hskip\z@skip
    \def\bbl@esphack{\loop\ifdim\lastskip>\z@\unskip\repeat\unskip}%
  \else
    \let\bbl@esphack\@empty
  \fi}
%    \end{macrocode}
%
% Ordered key/value lists (aka associative arrays). 
%
%    \begin{macrocode}
\def\bbl@array@end{\@empty}
\def\bbl@array@set#1#2#3{%
  \def\bbl@array@aux##1\bbl@elt#2\@@##2##3\bbl@array@aux##4{%
    \ifx\bbl@noelt##2%
      \def#1{##1\bbl@elt#2\@@{##4}}%
    \else
      \def\bbl@array@aux####1\bbl@elt#2\@@####2####3\bbl@array@aux####4{%
        \def#1{####1\bbl@elt#2\@@{####4}####3}}%
      \expandafter\bbl@array@aux#1\bbl@array@aux{#3}%
    \fi}%
  \expandafter\bbl@array@aux#1\bbl@elt#2\@@\bbl@noelt\bbl@array@aux{#3}}
\def\bbl@array@get#1#2#3{%
  \def\bbl@array@aux##1\bbl@elt#2\@@##2##3\bbl@array@aux{%
    \ifx\bbl@array@end##2\@empty\let#3\relax\else\def#3{##2}\fi}% 
  \expandafter\bbl@array@aux#1\bbl@elt#2\@@\bbl@array@end\bbl@array@aux}
%    \end{macrocode}
%
% That is all for the moment. Now follows some common stuff, for both
% Plain and \LaTeX{}. After it, we will resume the \LaTeX-only stuff.
%
% \section{Selectors and declaring languages}
%
% Still Plain + \texttt{base}.
%
%    \begin{macrocode}
\def\bbl@version{<@version@>}
\def\bbl@date{<@date@>}
<@Define Plain core switching macros@>
%    \end{macrocode}
%
% \macro{\adddialect}
% The macro |\adddialect| can be used to add the name of a dialect
% or variant language, for which an already defined hyphenation
% table can be used.
%
%    \begin{macrocode}
\def\adddialect#1#2{%
  \global\chardef#1#2\relax
  \bbl@usehooks{adddialect}{{#1}{#2}}%
  \begingroup
    \count@#1\relax
    \def\bbl@elt##1##2##3##4{%
      \ifnum\count@=##2\relax
        \edef\bbl@tempa{\expandafter\@gobbletwo\string#1}%
        \bbl@info{Hyphen rules for '\expandafter\@gobble\bbl@tempa'
                  set to \expandafter\string\csname l@##1\endcsname\\%
                  (\string\language\the\count@). Reported}%
        \def\bbl@elt####1####2####3####4{}%
      \fi}%
    \bbl@cs{languages}%
  \endgroup}
%    \end{macrocode}
%
% |\bbl@iflanguage| executes code only if the language |l@| exists.
%
%    \begin{macrocode}
\def\bbl@fixname#1{%
  \begingroup
    \def\bbl@tempe{l@}%
    \edef\bbl@tempd{\noexpand\@ifundefined{\noexpand\bbl@tempe#1}}%
    \bbl@tempd
      {\lowercase\expandafter{\bbl@tempd}%
         {\uppercase\expandafter{\bbl@tempd}%
           \@empty
           {\edef\bbl@tempd{\def\noexpand#1{#1}}%
            \uppercase\expandafter{\bbl@tempd}}}%
         {\edef\bbl@tempd{\def\noexpand#1{#1}}%
          \lowercase\expandafter{\bbl@tempd}}}%
      \@empty
    \edef\bbl@tempd{\endgroup\def\noexpand#1{#1}}%
  \bbl@tempd
  \bbl@exp{\\\bbl@usehooks{languagename}{{\languagename}{#1}}}}
\def\bbl@iflanguage#1{%
  \@ifundefined{l@#1}{\@nolanerr{#1}\@gobble}\@firstofone}
%    \end{macrocode}
%
% After a name has been ‘fixed’, the selectors will try to load the
% language. If even the fixed name is not defined, will load it on the
% fly, either based on its name, or if activated, its BCP 47 code.
%
%    \begin{macrocode}
\def\bbl@cntchars#1{\bbl@cntchars@i#1876543210\@@}
\def\bbl@cntchars@i#1#2#3#4#5#6#7#8#9\@@{\@car#9\@nil}%
\def\bbl@bcplookup#1{%
  \let\bbl@templ\@empty % language
  \let\bbl@temps\@empty % script
  \let\bbl@tempr\@empty % regions
  \let\bbl@tempv\@empty % variant
  \edef\bbl@tempa{#1}%
  \bbl@replace\bbl@tempa{-}{,}%
  \bbl@replace\bbl@tempa{_}{,}%
  \bbl@foreach\bbl@tempa{%
    \ifx\bbl@templ\@empty
      \lowercase{\def\bbl@templ{##1}}%
    \else\ifnum\bbl@cntchars{##1}=4
      \bbl@xin@{\@car##1\@nil}{0123456789}%
      \ifin@
        \def\bbl@tempv{##1}% eg, 1901
      \else
        \uppercase{\edef\bbl@tempb{\@car##1\@nil}}%
        \lowercase{\edef\bbl@temps{\bbl@tempb\@gobble##1}}%
      \fi
    \else\ifnum\bbl@cntchars{##1}<4
      \uppercase{\def\bbl@tempr{##1}}% 
    \else\ifnum\bbl@cntchars{##1}>4  
      \lowercase{\def\bbl@tempv{##1}}%
    \fi\fi\fi\fi}%
  \bbl@exp{\\\bbl@bcplookup@i{\bbl@templ}{\bbl@temps}{\bbl@tempr}{\bbl@tempv}}}
\def\bbl@bcplookup@try#1{%
 \ifx\bbl@bcp\relax
   \IfFileExists{babel-#1\bbl@tempe.ini}{\edef\bbl@bcp{#1\bbl@tempe}}{}%
\fi} 
\def\bbl@bcplookup@i#1#2#3#4{%
  \let\bbl@bcp\relax
  \def\bbl@tempe{#4}%
  \ifx\bbl@tempe\@empty\else
    \edef\bbl@tempe{-#4}%
  \fi
  \edef\bbl@tempa{#2#3\bbl@tempe}% en
  \ifx\bbl@tempa\@empty
    \bbl@bcplookup@try{#1}%
  \else
    \edef\bbl@tempa{#3\bbl@tempe}% en-Latn
    \ifx\bbl@tempa\@empty
      \bbl@bcplookup@try{#1-#2}%
      \bbl@bcplookup@try{#1}%
    \else
      \edef\bbl@tempa{#2\bbl@tempe}% en-US, en-001    
      \ifx\bbl@tempa\@empty
        \bbl@bcplookup@try{#1-#3}%
        \bbl@bcplookup@try{#1}%
      \else % en-Latn-US       
        \bbl@bcplookup@try{#1-#2-#3}%
        \bbl@bcplookup@try{#1-#2}%
        \bbl@bcplookup@try{#1-#3}%
        \bbl@bcplookup@try{#1}%
      \fi
    \fi
  \fi
  \ifx\bbl@bcp\relax
    \ifx\bbl@tempe\@empty\else\bbl@bcplookup@i{#1}{#2}{#3}{}\fi%
  \fi}
%
\let\bbl@initoload\relax
%    \end{macrocode}
%
% \macro{\iflanguage}
% Users might want to test (in a private package for instance)
% which language is currently active.
%
%    \begin{macrocode}
\def\iflanguage#1{%
  \bbl@iflanguage{#1}{%
    \ifnum\csname l@#1\endcsname=\language
      \expandafter\@firstoftwo
    \else
      \expandafter\@secondoftwo
    \fi}}
%    \end{macrocode}
%
% \subsection{Selecting the language}
%
% \macro{\selectlanguage}
%
% It checks whether the language is
% already defined before it performs its actual task, which is to
% update |\language| and activate language-specific definitions.
%
%    \begin{macrocode}
\let\bbl@select@type\z@
\edef\selectlanguage{%
  \noexpand\protect
  \expandafter\noexpand\csname selectlanguage \endcsname}
\ifx\@undefined\protect\let\protect\relax\fi
\let\xstring\string
%    \end{macrocode}
%
% \macro{\bbl@pop@language}
%
% \emph{But} when the language change happens \emph{inside} a group
% the end of the group doesn't write anything to the auxiliary
% files.
%
% \macro{\bbl@language@stack}
% The previous solution works for one level of nesting groups, but
% as soon as more levels are used it is no longer adequate.
%
%    \begin{macrocode}
\def\bbl@language@stack{}
%    \end{macrocode}
%
% \macro{\bbl@push@language}
% \macro{\bbl@pop@language}
% The stack is simply a list of languagenames, separated with a `+'
% sign; the push function can be simple:
%
%    \begin{macrocode}
\def\bbl@push@language{%
  \ifx\languagename\@undefined\else
    \ifx\currentgrouplevel\@undefined
      \xdef\bbl@language@stack{\languagename+\bbl@language@stack}%
    \else
      \ifnum\currentgrouplevel=\z@
        \xdef\bbl@language@stack{\languagename+}%
      \else
        \xdef\bbl@language@stack{\languagename+\bbl@language@stack}%
      \fi
    \fi
  \fi}
%    \end{macrocode}
%
% Retrieving information from the stack is a little bit less simple, as
% we need to remove the element from the stack while storing it in the
% macro |\languagename|.
%
% \macro{\bbl@pop@lang}
% This macro stores its first element (which is delimited by the
% `+'-sign) in |\languagename| and stores the rest of the string in
% |\bbl@language@stack|.
%
%    \begin{macrocode}
\def\bbl@pop@lang#1+#2\@@{%
  \edef\languagename{#1}%
  \xdef\bbl@language@stack{#2}}
\let\bbl@ifrestoring\@secondoftwo
\def\bbl@pop@language{%
  \expandafter\bbl@pop@lang\bbl@language@stack\@@
  \let\bbl@ifrestoring\@firstoftwo
  \expandafter\bbl@set@language\expandafter{\languagename}%
  \let\bbl@ifrestoring\@secondoftwo}
%    \end{macrocode}
%
% An alternative way to identify languages (in the \babel{} sense)
% with a numerical value is introduced in 3.30. 
%
%    \begin{macrocode}
\chardef\localeid\z@
\gdef\bbl@id@last{0}    % No real need for a new counter
\def\bbl@id@assign{%
  \bbl@ifunset{bbl@id@@\languagename}%
    {\count@\bbl@id@last\relax
     \advance\count@\@ne
     \global\bbl@csarg\chardef{id@@\languagename}\count@
     \xdef\bbl@id@last{\the\count@}%
     \ifcase\bbl@engine\or
       \directlua{
         Babel.locale_props[\bbl@id@last] = {}
         Babel.locale_props[\bbl@id@last].name = '\languagename'
         Babel.locale_props[\bbl@id@last].vars = {}
        }%
      \fi}%
    {}%
    \chardef\localeid\bbl@cl{id@}}
%    \end{macrocode}
%
% The unprotected part of |\selectlanguage|. In case it is used as
% environment, declare |\endselectlaguage|, just for safety.
%
%    \begin{macrocode}
\let\bbl@select@opts\@empty
\expandafter\def\csname selectlanguage \endcsname{%
  \@ifnextchar[\bbl@select@s{\bbl@select@s[]}}
\def\bbl@select@s[#1]#2{%
  \bbl@forkv{#1}{\bbl@array@set\bbl@select@opts{##1}{##2}}%
  \ifnum\bbl@hymapsel=\@cclv\let\bbl@hymapsel\tw@\fi
  \bbl@push@language
  \aftergroup\bbl@pop@language
  \bbl@set@language{#2}}
\let\endselectlanguage\relax
%    \end{macrocode}
%
% \macro{\bbl@set@language}
%
% The macro |\bbl@set@language| takes care of switching the language
% environment \emph{and} of writing entries on the auxiliary files.
%
%    \begin{macrocode}
\def\BabelContentsFiles{toc,lof,lot}
\def\bbl@set@language#1{% from selectlanguage, pop@
  % The old buggy way. Preserved for compatibility, but simplified
  \edef\languagename{\expandafter\string#1\@empty}%
  \select@language{\languagename}%
  \bbl@array@get\bbl@select@opts{main}\bbl@tempa
  \ifx\bbl@tempa\relax\else
    \let\bbl@main@language\localename
    \let\mainlocalename\localename
  \fi
  % write to aux files
  \expandafter\ifx\csname date\languagename\endcsname\relax\else
    \if@filesw
      \bbl@array@get\bbl@select@opts{nofiles}\bbl@tempa
      \ifx\bbl@tempa\relax
        \ifx\babel@aux\@gobbletwo\else % Set if single in the first, redundant
          \bbl@savelastskip
          \protected@write\@auxout{}{\string\babel@aux{\bbl@auxname}{}}%
          \bbl@restorelastskip
        \fi
        \bbl@usehooks{write}{}%
      \fi
    \fi
  \fi
  \let\bbl@select@opts\@empty}
%
\let\bbl@restorelastskip\relax
\let\bbl@savelastskip\relax
%
\def\select@language#1{% from set@, babel@aux, babel@toc
  \ifx\bbl@selectorname\@empty
    \def\bbl@selectorname{select}%
  \fi
  % set hymap
  \ifnum\bbl@hymapsel=\@cclv\chardef\bbl@hymapsel4\relax\fi
  % set name (when coming from babel@aux)
  \edef\languagename{#1}%
  \bbl@fixname\languagename
  % define \localename when coming from set@, with a trick 
  \ifx\scantokens\@undefined
    \def\localename{??}%
  \else
    \bbl@exp{\\\scantokens{\def\\\localename{\languagename}\\\noexpand}\relax}%
  \fi
  \bbl@provide@locale
  \bbl@iflanguage\languagename{%
    \let\bbl@select@type\z@
    \expandafter\bbl@switch\expandafter{\languagename}}}
\def\babel@aux#1#2{%
  \select@language{#1}%
  \bbl@foreach\BabelContentsFiles{% \relax -> don't assume vertical mode
    \@writefile{##1}{\babel@toc{#1}{#2}\relax}}}%
\def\babel@toc#1#2{%
  \select@language{#1}}
%    \end{macrocode}
%
% First, check if the user asks for a known language. If so,
% update the value of |\language| and call |\originalTeX|
% to bring \TeX\ in a certain pre-defined state.
%
%    \begin{macrocode}
\newif\ifbbl@usedategroup
\let\bbl@savedextras\@empty
\def\bbl@switch#1{%  from select@, foreign@
  % restore
  \originalTeX
  \expandafter\def\expandafter\originalTeX\expandafter{%
    \csname noextras#1\endcsname
    \let\originalTeX\@empty
    \babel@beginsave}%
  \bbl@usehooks{afterreset}{}%
  \languageshorthands{none}%
  % set the locale id
  \bbl@id@assign
  % switch captions, date
  \bbl@bsphack
    \ifcase\bbl@select@type
      \csname captions#1\endcsname\relax
      \csname date#1\endcsname\relax
    \else
      \bbl@array@get\bbl@select@opts{captions}\bbl@tempa
      \ifx\bbl@tempa\relax\else
        \csname captions#1\endcsname\relax
      \fi
      \bbl@array@get\bbl@select@opts{date}\bbl@tempa
      \ifx\bbl@tempa\relax\else % if \foreign... within \<language>date
        \csname date#1\endcsname\relax
      \fi
    \fi
  \bbl@esphack
  % switch extras
  \csname bbl@preextras@#1\endcsname
  \bbl@usehooks{beforeextras}{}%
  \csname extras#1\endcsname\relax
  \bbl@usehooks{afterextras}{}%
  %  > babel-ensure
  %  > babel-sh-<short>
  %  > babel-bidi
  %  > babel-fontspec
  \let\bbl@savedextras\@empty
  % hyphenation - case mapping
  \ifcase\bbl@opt@hyphenmap\or
    \def\BabelLower##1##2{\lccode##1=##2\relax}%
    \ifnum\bbl@hymapsel>4\else
      \csname\languagename @bbl@hyphenmap\endcsname
    \fi
    \chardef\bbl@opt@hyphenmap\z@
  \else
    \ifnum\bbl@hymapsel>\bbl@opt@hyphenmap\else
      \csname\languagename @bbl@hyphenmap\endcsname
    \fi
  \fi
  \let\bbl@hymapsel\@cclv
  % hyphenation - select rules
  \ifnum\csname l@\languagename\endcsname=\l@unhyphenated
    \edef\bbl@tempa{u}%
  \else
    \edef\bbl@tempa{\bbl@cl{lnbrk}}%
  \fi
  % linebreaking - handle u, e, k (v in the future)
  \bbl@xin@{/u}{/\bbl@tempa}%
  \ifin@\else\bbl@xin@{/e}{/\bbl@tempa}\fi % elongated forms 
  \ifin@\else\bbl@xin@{/k}{/\bbl@tempa}\fi % only kashida
  \ifin@\else\bbl@xin@{/p}{/\bbl@tempa}\fi % padding (e.g., Tibetan)
  \ifin@\else\bbl@xin@{/v}{/\bbl@tempa}\fi % variable font
  % hyphenation - save mins
  \babel@savevariable\lefthyphenmin
  \babel@savevariable\righthyphenmin
  \ifnum\bbl@engine=\@ne
    \babel@savevariable\hyphenationmin
  \fi
  \ifin@
    % unhyphenated/kashida/elongated/padding = allow stretching
    \language\l@unhyphenated
    \babel@savevariable\emergencystretch
    \emergencystretch\maxdimen
    \babel@savevariable\hbadness
    \hbadness\@M
  \else
    % other = select patterns
    \bbl@patterns{#1}%
  \fi
  % hyphenation - set mins
  \expandafter\ifx\csname #1hyphenmins\endcsname\relax
    \set@hyphenmins\tw@\thr@@\relax
    \@nameuse{bbl@hyphenmins@}%
  \else
    \expandafter\expandafter\expandafter\set@hyphenmins
      \csname #1hyphenmins\endcsname\relax
  \fi
  \@nameuse{bbl@hyphenmins@}%
  \@nameuse{bbl@hyphenmins@\languagename}%
  \@nameuse{bbl@hyphenatmin@}%
  \@nameuse{bbl@hyphenatmin@\languagename}%
  \let\bbl@selectorname\@empty}
%    \end{macrocode}
%
% \environment{otherlanguage}
% It can be used as an alternative to
% using the |\selectlanguage| declarative command. 
%
%    \begin{macrocode}
\edef\otherlanguage{%
  \noexpand\protect
  \expandafter\noexpand\csname otherlanguage \endcsname}
\expandafter\def\csname otherlanguage \endcsname{%
  \@ifstar{\@nameuse{otherlanguage*}}\bbl@otherlanguage}
\def\bbl@otherlanguage#1{%
  \def\bbl@selectorname{other}%
  \ifnum\bbl@hymapsel=\@cclv\let\bbl@hymapsel\thr@@\fi
  \csname selectlanguage \endcsname{#1}%
  \ignorespaces}
%    \end{macrocode}
%
% The |\endotherlanguage| part of the environment tries to hide
% itself when it is called in horizontal mode.
%
%    \begin{macrocode}
\long\def\endotherlanguage{\@ignoretrue\ignorespaces}
%    \end{macrocode}
%
% \environment{otherlanguage*}
% It is meant to be used when a
% large part of text from a different language needs to be typeset,
% but without changing the translation of words such as `figure'.
%
%    \begin{macrocode}
\expandafter\def\csname otherlanguage*\endcsname{%
  \@ifnextchar[\bbl@otherlanguage@s{\bbl@otherlanguage@s[]}}
\def\bbl@otherlanguage@s[#1]#2{%
  \def\bbl@selectorname{other*}%
  \ifnum\bbl@hymapsel=\@cclv\chardef\bbl@hymapsel4\relax\fi
  \bbl@forkv{#1}{\bbl@array@set\bbl@select@opts{##1}{##2}}%
  \foreign@language{#2}}
\expandafter\let\csname endotherlanguage*\endcsname\relax
%    \end{macrocode}
%
% \macro{\foreignlanguage}
%
% This command takes two arguments, the first argument is the name of
% the language to use for typesetting the text specified in the second
% argument.
%
%    \begin{macrocode}
\providecommand\bbl@beforeforeign{}
\edef\foreignlanguage{%
  \noexpand\protect
  \expandafter\noexpand\csname foreignlanguage \endcsname}
\expandafter\def\csname foreignlanguage \endcsname{%
  \@ifstar\bbl@foreign@s\bbl@foreign@x}
\providecommand\bbl@foreign@x[3][]{%
  \begingroup
    \def\bbl@selectorname{foreign}%
    \bbl@forkv{#1}{\bbl@array@set\bbl@select@opts{##1}{##2}}%
    \let\BabelText\@firstofone
    \bbl@beforeforeign
    \foreign@language{#2}%
    \bbl@usehooks{foreign}{}%
    \BabelText{#3}% Now in horizontal mode!
  \endgroup}
\def\bbl@foreign@s#1#2{%
  \begingroup
    {\par}%
    \def\bbl@selectorname{foreign*}%
    \let\bbl@select@opts\@empty
    \let\BabelText\@firstofone
    \foreign@language{#1}%
    \bbl@usehooks{foreign*}{}%
    \bbl@dirparastext
    \BabelText{#2}% Still in vertical mode!
    {\par}%
  \endgroup}
\providecommand\BabelWrapText[1]{%
   \def\bbl@tempa{\def\BabelText####1}%
   \expandafter\bbl@tempa\expandafter{\BabelText{#1}}}
%    \end{macrocode}
%
% \macro{\foreign@language}
%
% This macro does the work for |\foreignlanguage| and the
% \Lenv{otherlanguage*} environment. 
%
%    \begin{macrocode}
\def\foreign@language#1{%
  % set name
  \edef\languagename{#1}%
  \ifbbl@usedategroup
    \bbl@array@set\bbl@select@opts{date}{\@empty}%
    \bbl@usedategroupfalse
  \fi
  \bbl@fixname\languagename
  \let\localename\languagename
  \bbl@provide@locale
  \bbl@iflanguage\languagename{%
    \let\bbl@select@type\@ne
    \expandafter\bbl@switch\expandafter{\languagename}}}
%    \end{macrocode}
%
% The following macro executes conditionally some code based on the
% selector being used.
%
%    \begin{macrocode}
\def\IfBabelSelectorTF#1{%
  \bbl@xin@{,\bbl@selectorname,}{,\zap@space#1 \@empty,}%
  \ifin@
    \expandafter\@firstoftwo
  \else
    \expandafter\@secondoftwo
  \fi}
%    \end{macrocode}
%
% \macro{\bbl@patterns}
%
% This macro selects the hyphenation patterns by changing the
% \cs{language} register. 
%
%    \begin{macrocode}
\let\bbl@hyphlist\@empty
\let\bbl@hyphenation@\relax
\let\bbl@pttnlist\@empty
\let\bbl@patterns@\relax
\let\bbl@hymapsel=\@cclv
\def\bbl@patterns#1{%
  \language=\expandafter\ifx\csname l@#1:\f@encoding\endcsname\relax
      \csname l@#1\endcsname
      \edef\bbl@tempa{#1}%
    \else
      \csname l@#1:\f@encoding\endcsname
      \edef\bbl@tempa{#1:\f@encoding}%
    \fi
  \@expandtwoargs\bbl@usehooks{patterns}{{#1}{\bbl@tempa}}%
  %  > luatex
  \@ifundefined{bbl@hyphenation@}{}{% Can be \relax!
    \begingroup
      \bbl@xin@{,\number\language,}{,\bbl@hyphlist}%
      \ifin@\else
        \@expandtwoargs\bbl@usehooks{hyphenation}{{#1}{\bbl@tempa}}%
        \hyphenation{%
          \bbl@hyphenation@
          \@ifundefined{bbl@hyphenation@#1}%
            \@empty
            {\space\csname bbl@hyphenation@#1\endcsname}}%
        \xdef\bbl@hyphlist{\bbl@hyphlist\number\language,}%
      \fi
    \endgroup}}
%    \end{macrocode}
%
% \environment{hyphenrules}
%
% It can be used to select \emph{just} the hyphenation rules.
%
%    \begin{macrocode}
\def\hyphenrules#1{%
  \edef\bbl@tempf{#1}%
  \bbl@fixname\bbl@tempf
  \bbl@iflanguage\bbl@tempf{%
    \expandafter\bbl@patterns\expandafter{\bbl@tempf}%
    \ifx\languageshorthands\@undefined\else
      \languageshorthands{none}%
    \fi
    \expandafter\ifx\csname\bbl@tempf hyphenmins\endcsname\relax
      \set@hyphenmins\tw@\thr@@\relax
    \else
      \expandafter\expandafter\expandafter\set@hyphenmins
      \csname\bbl@tempf hyphenmins\endcsname\relax
    \fi}}
\let\endhyphenrules\@empty
%    \end{macrocode}
%
% \macro{\providehyphenmins}
% The macro |\providehyphenmins| should be used in the language
% definition files to provide a \emph{default} setting for the
% hyphenation parameters |\lefthyphenmin| and |\righthyphenmin|
%
%    \begin{macrocode}
\def\providehyphenmins#1#2{%
  \expandafter\ifx\csname #1hyphenmins\endcsname\relax
    \@namedef{#1hyphenmins}{#2}%
  \fi}
%    \end{macrocode}
%
% \macro{\set@hyphenmins}
% This macro sets the values of |\lefthyphenmin| and |\righthyphenmin|.
% It expects two values as its argument.
%
%    \begin{macrocode}
\def\set@hyphenmins#1#2{%
  \lefthyphenmin#1\relax
  \righthyphenmin#2\relax}
%    \end{macrocode}
%
% \macro{\ProvidesLanguage}
% The identification code for each file is something that was
% introduced in \LaTeXe. 
%
%    \begin{macrocode}
\ifx\ProvidesFile\@undefined
  \def\ProvidesLanguage#1[#2 #3 #4]{%
    \wlog{Language: #1 #4 #3 <#2>}%
    }
\else
  \def\ProvidesLanguage#1{%
    \begingroup
      \catcode`\ 10 %
      \@makeother\/%
      \@ifnextchar[%]
        {\@provideslanguage{#1}}{\@provideslanguage{#1}[]}}
  \def\@provideslanguage#1[#2]{%
    \wlog{Language: #1 #2}%
    \expandafter\xdef\csname ver@#1.ldf\endcsname{#2}%
    \endgroup}
\fi
\ifx\originalTeX\@undefined\let\originalTeX\@empty\fi
\ifx\babel@beginsave\@undefined\let\babel@beginsave\relax\fi
%    \end{macrocode}
%
% A few macro names are reserved for future releases of \babel, which
% will use the concept of `locale':
%
%    \begin{macrocode}
\providecommand\setlocale{\bbl@error{not-yet-available}{}{}{}}
\let\uselocale\setlocale
\let\locale\setlocale
\let\selectlocale\setlocale
\let\textlocale\setlocale
\let\textlanguage\setlocale
\let\languagetext\setlocale
%    \end{macrocode}
%
% \subsection{Errors}
%
% \macro{\@nolanerr}
% \macro{\@nopatterns}
%
% The \babel\ package will signal an error when a documents tries
% to select a language that hasn't been defined earlier. 
%
% \macro{\@noopterr}
% When the package was loaded without options not everything will
% work as expected. An error message is issued in that case.
%
%    \begin{macrocode}
\edef\bbl@nulllanguage{\string\language=0}
\def\bbl@nocaption#1#2{\NoCaseChange{\protect\bbl@nocaption@i{#1}{#2}}}
\def\bbl@nocaption@i#1#2{% 1: text to be printed 2: caption macro \langXname
  \global\@namedef{#2}{\textbf{?#1?}}%
  \@nameuse{#2}%
  \edef\bbl@tempa{#1}%
  \bbl@sreplace\bbl@tempa{name}{}%
  \bbl@warning{%
    \@backslashchar#1 not set for '\languagename'. Please,\\%
    define it after the language has been loaded\\%
    (typically in the preamble) with:\\%
    \string\setlocalecaption{\languagename}{\bbl@tempa}{..}\\%
    Feel free to contribute on github.com/latex3/babel.\\%
    Reported}}
\def\bbl@tentative{\protect\bbl@tentative@i}
\def\bbl@tentative@i#1{%
  \bbl@warning{%
    Some functions for '#1' are tentative.\\%
    They might not work as expected and their behavior\\%
    could change in the future.\\%
    Reported}}
\def\@nolanerr#1{\bbl@error{undefined-language}{#1}{}{}}
\def\@nopatterns#1{%
  \bbl@warning
    {No hyphenation patterns were preloaded for\\%
     the language '#1' into the format.\\%
     Please, configure your TeX system to add them and\\%
     rebuild the format. Now I will use the patterns\\%
     preloaded for \bbl@nulllanguage\space instead}}
\let\bbl@usehooks\@gobbletwo
%    \end{macrocode}
%
% \section{Finish the very basic macros}
% \label{endswitch}
%
% Here ended the \textsf{base} option (\LaTeX).
%
% \section{Only Plain}
%
%    \begin{macrocode}
\ifx\bbl@onlyswitch\@empty\endinput\fi
%    \end{macrocode}
%
% \subsection{More on selection}
%
% \macro{\babelensure}
%
% The user command just parses the optional argument and creates a
% new macro named |\bbl@e@<language>|. 
%
%    \begin{macrocode}
\bbl@trace{Defining babelensure}
\newcommand\babelensure[2][]{%
  \AddBabelHook{babel-ensure}{afterextras}{%
    \ifcase\bbl@select@type
      \bbl@cl{e}%
    \fi}%
  \begingroup
    \let\bbl@ens@include\@empty
    \let\bbl@ens@exclude\@empty
    \def\bbl@ens@fontenc{\relax}%
    \def\bbl@tempb##1{%
      \ifx\@empty##1\else\noexpand##1\expandafter\bbl@tempb\fi}%
    \edef\bbl@tempa{\bbl@tempb#1\@empty}%
    \def\bbl@tempb##1=##2\@@{\@namedef{bbl@ens@##1}{##2}}%
    \bbl@foreach\bbl@tempa{\bbl@tempb##1\@@}%
    \def\bbl@tempc{\bbl@ensure}%
    \expandafter\bbl@add\expandafter\bbl@tempc\expandafter{%
      \expandafter{\bbl@ens@include}}%
    \expandafter\bbl@add\expandafter\bbl@tempc\expandafter{%
      \expandafter{\bbl@ens@exclude}}%
    \toks@\expandafter{\bbl@tempc}%
    \bbl@exp{%
  \endgroup
  \def\<bbl@e@#2>{\the\toks@{\bbl@ens@fontenc}}}}
\def\bbl@ensure#1#2#3{% 1: include 2: exclude 3: fontenc
  \def\bbl@tempb##1{% elt for (excluding) \bbl@captionslist list
    \ifx##1\@undefined % 3.32 - Don't assume the macro exists
      \edef##1{\noexpand\bbl@nocaption
        {\bbl@stripslash##1}{\languagename\bbl@stripslash##1}}%
    \fi
    \ifx##1\@empty\else
      \in@{##1}{#2}%
      \ifin@\else
        \let\bbl@tempc\languagename
        \bbl@replace\bbl@tempc{-}{@@}%
        \bbl@ifunset{bbl@ensure@\bbl@tempc}%
          {\bbl@exp{%
            \\\DeclareRobustCommand\<bbl@ensure@\bbl@tempc>[1]{%
              \\\foreignlanguage{\languagename}%
              {\ifx\relax#3\else
                \\\fontencoding{#3}\\\selectfont
               \fi
               ########1}}}}%
          {}%
        \toks@\expandafter{##1}%
        \edef##1{%
           \bbl@csarg\noexpand{ensure@\bbl@tempc}%
           {\the\toks@}}%
      \fi
      \expandafter\bbl@tempb
    \fi}%
  \expandafter\bbl@tempb\bbl@captionslist\today\@empty
  \def\bbl@tempa##1{% elt for include list
    \ifx##1\@empty\else
      \let\bbl@tempc\languagename
      \bbl@replace\bbl@tempc{-}{@@}%
      \bbl@csarg\in@{ensure@\bbl@tempc\expandafter}\expandafter{##1}%
      \ifin@\else
        \bbl@tempb##1\@empty
      \fi
      \expandafter\bbl@tempa
    \fi}%
  \bbl@tempa#1\@empty}
\def\bbl@captionslist{%
  \prefacename\refname\abstractname\bibname\chaptername\appendixname
  \contentsname\listfigurename\listtablename\indexname\figurename
  \tablename\partname\enclname\ccname\headtoname\pagename\seename
  \alsoname\proofname\glossaryname}
%    \end{macrocode}
%
% \subsection{Short tags}
%
% \macro{\babeltags}
% This macro is straightforward. 
%
%   \begin{macrocode}
\bbl@trace{Short tags}
\newcommand\babeltags[1]{%
  \edef\bbl@tempa{\zap@space#1 \@empty}%
  \def\bbl@tempb##1=##2\@@{%
    \edef\bbl@tempc{%
      \noexpand\newcommand
      \expandafter\noexpand\csname ##1\endcsname{%
        \noexpand\protect
        \expandafter\noexpand\csname otherlanguage*\endcsname{##2}}
      \noexpand\newcommand
      \expandafter\noexpand\csname text##1\endcsname{%
        \noexpand\foreignlanguage{##2}}}
    \bbl@tempc}%
  \bbl@for\bbl@tempa\bbl@tempa{%
    \expandafter\bbl@tempb\bbl@tempa\@@}}
%    \end{macrocode}
%
% \subsection{Compatibility with language.def}
%
% Plain e-\TeX{} doesn’t rely on language.dat, but babel can be made
% compatible with this format easily.
%
%    \begin{macrocode}
\bbl@trace{Compatibility with language.def}
\ifx\directlua\@undefined\else
  \ifx\bbl@luapatterns\@undefined
    \input luababel.def
  \fi
\fi
\ifx\bbl@languages\@undefined
  \ifx\directlua\@undefined
    \openin0 = language.def
    \ifeof0
      \closein0
      \message{I couldn't find the file language.def}
    \else
      \closein0
      \begingroup
        \def\addlanguage#1#2#3#4#5{%
          \expandafter\ifx\csname lang@#1\endcsname\relax\else
            \global\expandafter\let\csname l@#1\expandafter\endcsname
              \csname lang@#1\endcsname
          \fi}%
        \def\uselanguage#1{}%
        \input language.def
      \endgroup
    \fi
  \fi
  \chardef\l@english\z@
\fi
%    \end{macrocode}
%
% \macro{\addto}
% It takes two arguments, a \meta{control sequence} and \TeX-code to
% be added to the \meta{control sequence}.
%
%    \begin{macrocode}
\def\addto#1#2{%
  \ifx#1\@undefined
    \def#1{#2}%
  \else
    \ifx#1\relax
      \def#1{#2}%
    \else
      {\toks@\expandafter{#1#2}%
       \xdef#1{\the\toks@}}%
    \fi
  \fi}
%    \end{macrocode}
%
% \subsection{Hooks}
%
% Admittedly, the current implementation is a somewhat simplistic.
%
%    \begin{macrocode}
\bbl@trace{Hooks}
\newcommand\AddBabelHook[3][]{%
  \bbl@ifunset{bbl@hk@#2}{\EnableBabelHook{#2}}{}%
  \def\bbl@tempa##1,#3=##2,##3\@empty{\def\bbl@tempb{##2}}%
  \expandafter\bbl@tempa\bbl@evargs,#3=,\@empty
  \bbl@ifunset{bbl@ev@#2@#3@#1}%
    {\bbl@csarg\bbl@add{ev@#3@#1}{\bbl@elth{#2}}}%
    {\bbl@csarg\let{ev@#2@#3@#1}\relax}%
  \bbl@csarg\newcommand{ev@#2@#3@#1}[\bbl@tempb]}
\newcommand\EnableBabelHook[1]{\bbl@csarg\let{hk@#1}\@firstofone}
\newcommand\DisableBabelHook[1]{\bbl@csarg\let{hk@#1}\@gobble}
\def\bbl@usehooks{\bbl@usehooks@lang\languagename}
\def\bbl@usehooks@lang#1#2#3{% Test for Plain
  \ifx\UseHook\@undefined\else\UseHook{babel/*/#2}\fi
  \def\bbl@elth##1{%
    \bbl@cs{hk@##1}{\bbl@cs{ev@##1@#2@}#3}}%
  \bbl@cs{ev@#2@}%
  \ifx\languagename\@undefined\else % Test required for Plain (?)
    \ifx\UseHook\@undefined\else\UseHook{babel/#1/#2}\fi
    \def\bbl@elth##1{%
      \bbl@cs{hk@##1}{\bbl@cs{ev@##1@#2@#1}#3}}%
    \bbl@cs{ev@#2@#1}%
  \fi}
\def\bbl@evargs{,% <- don't delete this comma
  everylanguage=1,loadkernel=1,loadpatterns=1,loadexceptions=1,%
  adddialect=2,patterns=2,defaultcommands=0,encodedcommands=2,write=0,%
  beforeextras=0,afterextras=0,stopcommands=0,stringprocess=0,%
  hyphenation=2,initiateactive=3,afterreset=0,foreign=0,foreign*=0,%
  beforestart=0,languagename=2,begindocument=1}
\ifx\NewHook\@undefined\else % Test for Plain (?)
  \def\bbl@tempa#1=#2\@@{\NewHook{babel/#1}\NewHook{babel/*/#1}}
  \bbl@foreach\bbl@evargs{\bbl@tempa#1\@@}
\fi
%    \end{macrocode}
%
% \subsection{Setting up language files}
%
% \macro{\LdfInit}
%
% |\LdfInit| macro takes two arguments. The first
% argument is the name of the language that will be defined in the
% language definition file; the second argument is either a control
% sequence or a string from which a control sequence should be
% constructed.
%
%    \begin{macrocode}
\bbl@trace{Macros for setting language files up}
\def\bbl@ldfinit{%
  \let\bbl@screset\@empty
  \let\BabelStrings\bbl@opt@string
  \let\BabelOptions\@empty
  \let\BabelLanguages\relax
  \ifx\originalTeX\@undefined
    \let\originalTeX\@empty
  \else
    \originalTeX
  \fi}
\def\LdfInit#1#2{%
  \chardef\atcatcode=\catcode`\@
  \catcode`\@=11\relax
  \chardef\eqcatcode=\catcode`\=
  \catcode`\==12\relax
  \@ifpackagewith{babel}{ensureinfo=off}{}%
    {\ifx\InputIfFileExists\@undefined\else
       \bbl@ifunset{bbl@lname@#1}%
         {{\let\bbl@ensuring\@empty % Flag used in babel-serbianc.tex
          \def\languagename{#1}%
          \bbl@id@assign
          \bbl@load@info{#1}}}%
        {}%
     \fi}%
  \expandafter\if\expandafter\@backslashchar
                 \expandafter\@car\string#2\@nil
    \ifx#2\@undefined\else
      \ldf@quit{#1}%
    \fi
  \else
    \expandafter\ifx\csname#2\endcsname\relax\else
      \ldf@quit{#1}%
    \fi
  \fi
  \bbl@ldfinit}
%    \end{macrocode}
%
% \macro{\ldf@quit}
% This macro interrupts the processing of a language definition file.
%
%    \begin{macrocode}
\def\ldf@quit#1{%
  \expandafter\main@language\expandafter{#1}%
  \catcode`\@=\atcatcode \let\atcatcode\relax
  \catcode`\==\eqcatcode \let\eqcatcode\relax
  \endinput}
%    \end{macrocode}
%
% \macro{\ldf@finish}
% This macro takes one argument. It is the name of the language that
% was defined in the language definition file.
%
%    \begin{macrocode}
\def\bbl@afterldf{%
  \bbl@afterlang
  \let\bbl@afterlang\relax
  \let\BabelModifiers\relax
  \let\bbl@screset\relax}%
\def\ldf@finish#1{%
  \loadlocalcfg{#1}%
  \bbl@afterldf
  \expandafter\main@language\expandafter{#1}%
  \catcode`\@=\atcatcode \let\atcatcode\relax
  \catcode`\==\eqcatcode \let\eqcatcode\relax}
%    \end{macrocode}
%
% After the preamble of the document the commands |\LdfInit|,
% |\ldf@quit| and |\ldf@finish| are no longer needed. 
%
%    \begin{macrocode}
\@onlypreamble\LdfInit
\@onlypreamble\ldf@quit
\@onlypreamble\ldf@finish
%    \end{macrocode}
%
% \macro{\main@language}
% \macro{\bbl@main@language}
% This command should be used in the various language definition files.
%
%    \begin{macrocode}
\def\main@language#1{%
  \def\bbl@main@language{#1}%
  \let\languagename\bbl@main@language
  \let\localename\bbl@main@language
  \let\mainlocalename\bbl@main@language
  \bbl@id@assign
  \ifcase\bbl@engine\or
    \ifx\setattribute\@undefined\else
      \setattribute\bbl@attr@locale\localeid
    \fi
  \fi
  \bbl@patterns{\languagename}}
%    \end{macrocode}
%
% We also have to make sure that some code gets executed at the
% beginning of the document.
%
%    \begin{macrocode}
\AtBeginDocument{\expandafter\selectlanguage\expandafter{\bbl@main@language}}
%    \end{macrocode}
%
% A bit of optimization. Select in heads/feet the language only if
% necessary.
%
%    \begin{macrocode}
\def\select@language@x#1{%
  \ifcase\bbl@select@type
    \bbl@ifsamestring\languagename{#1}{}{\select@language{#1}}%
  \else
    \select@language{#1}%
  \fi}
%    \end{macrocode}
%
% \subsection{Shorthands}
%
% The macro |\initiate@active@char| below takes all the necessary
% actions to make its argument a shorthand character.
%
%    \begin{macrocode}
\bbl@trace{Shorhands}
\def\bbl@withactive#1#2{%
  \begingroup
    \lccode`~=`#2\relax
    \lowercase{\endgroup#1~}}
%    \end{macrocode}
%
% \macro{\bbl@add@special}
% The macro |\bbl@add@special| is used to add a new character (or
% single character control sequence) to the macro |\dospecials| (and
% |\@sanitize| if \LaTeX\ is used).
%
%    \begin{macrocode}
\def\bbl@add@special#1{% 1:a macro like \", \?, etc.
  \bbl@add\dospecials{\do#1}% test @sanitize = \relax, for back. compat.
  \bbl@ifunset{@sanitize}{}{\bbl@add\@sanitize{\@makeother#1}}%
  \ifx\nfss@catcodes\@undefined\else
    \begingroup
      \catcode`#1\active
      \nfss@catcodes
      \ifnum\catcode`#1=\active
        \endgroup
        \bbl@add\nfss@catcodes{\@makeother#1}%
      \else
        \endgroup
      \fi
  \fi}
%    \end{macrocode}
%
% \macro{\initiate@active@char}
% A language definition file can call this macro to make a character
% active. 
%
%    \begin{macrocode}
\def\bbl@active@def#1#2#3#4{%
  \@namedef{#3#1}{%
    \expandafter\ifx\csname#2@sh@#1@\endcsname\relax
      \bbl@afterelse\bbl@sh@select#2#1{#3@arg#1}{#4#1}%
    \else
      \bbl@afterfi\csname#2@sh@#1@\endcsname
    \fi}%
  \long\@namedef{#3@arg#1}##1{%
    \expandafter\ifx\csname#2@sh@#1@\string##1@\endcsname\relax
      \bbl@afterelse\csname#4#1\endcsname##1%
    \else
      \bbl@afterfi\csname#2@sh@#1@\string##1@\endcsname
    \fi}}%
\def\initiate@active@char#1{%
  \bbl@ifunset{active@char\string#1}%
    {\bbl@withactive
      {\expandafter\@initiate@active@char\expandafter}#1\string#1#1}%
    {}}
\def\@initiate@active@char#1#2#3{%
  \bbl@csarg\edef{oricat@#2}{\catcode`#2=\the\catcode`#2\relax}%
  \ifx#1\@undefined
    \bbl@csarg\def{oridef@#2}{\def#1{\active@prefix#1\@undefined}}%
  \else
    \bbl@csarg\let{oridef@@#2}#1%
    \bbl@csarg\edef{oridef@#2}{%
      \let\noexpand#1%
      \expandafter\noexpand\csname bbl@oridef@@#2\endcsname}%
  \fi
  \ifx#1#3\relax
    \expandafter\let\csname normal@char#2\endcsname#3%
  \else
    \bbl@info{Making #2 an active character}%
    \ifnum\mathcode`#2=\ifodd\bbl@engine"1000000 \else"8000 \fi
      \@namedef{normal@char#2}{%
        \textormath{#3}{\csname bbl@oridef@@#2\endcsname}}%
    \else
      \@namedef{normal@char#2}{#3}%
    \fi
    \bbl@restoreactive{#2}%
    \AtBeginDocument{%
      \catcode`#2\active
      \if@filesw
        \immediate\write\@mainaux{\catcode`\string#2\active}%
      \fi}%
    \expandafter\bbl@add@special\csname#2\endcsname
    \catcode`#2\active
  \fi
  \let\bbl@tempa\@firstoftwo
  \if\string^#2%
    \def\bbl@tempa{\noexpand\textormath}%
  \else
    \ifx\bbl@mathnormal\@undefined\else
      \let\bbl@tempa\bbl@mathnormal
    \fi
  \fi
  \expandafter\edef\csname active@char#2\endcsname{%
    \bbl@tempa
      {\noexpand\if@safe@actives
         \noexpand\expandafter
         \expandafter\noexpand\csname normal@char#2\endcsname
       \noexpand\else
         \noexpand\expandafter
         \expandafter\noexpand\csname bbl@doactive#2\endcsname
       \noexpand\fi}%
     {\expandafter\noexpand\csname normal@char#2\endcsname}}%
  \bbl@csarg\edef{doactive#2}{%
    \expandafter\noexpand\csname user@active#2\endcsname}%
  \bbl@csarg\edef{active@#2}{%
    \noexpand\active@prefix\noexpand#1%
    \expandafter\noexpand\csname active@char#2\endcsname}%
  \bbl@csarg\edef{normal@#2}{%
    \noexpand\active@prefix\noexpand#1%
    \expandafter\noexpand\csname normal@char#2\endcsname}%
  \bbl@ncarg\let#1{bbl@normal@#2}%
  \bbl@active@def#2\user@group{user@active}{language@active}%
  \bbl@active@def#2\language@group{language@active}{system@active}%
  \bbl@active@def#2\system@group{system@active}{normal@char}%
  \expandafter\edef\csname\user@group @sh@#2@@\endcsname
    {\expandafter\noexpand\csname normal@char#2\endcsname}%
  \expandafter\edef\csname\user@group @sh@#2@\string\protect@\endcsname
    {\expandafter\noexpand\csname user@active#2\endcsname}%
  \if\string'#2%
    \let\prim@s\bbl@prim@s
    \let\active@math@prime#1%
  \fi
  \bbl@usehooks{initiateactive}{{#1}{#2}{#3}}}
\@ifpackagewith{babel}{KeepShorthandsActive}%
  {\let\bbl@restoreactive\@gobble}%
  {\def\bbl@restoreactive#1{%
     \bbl@exp{%
       \\\AfterBabelLanguage\\\CurrentOption
         {\catcode`#1=\the\catcode`#1\relax}%
       \\\AtEndOfPackage
         {\catcode`#1=\the\catcode`#1\relax}}}%
   \AtEndOfPackage{\let\bbl@restoreactive\@gobble}}
%    \end{macrocode}
%
% \macro{\bbl@sh@select}
% This command helps the shorthand supporting macros to select how to
% proceed.
%
%    \begin{macrocode}
\def\bbl@sh@select#1#2{%
  \expandafter\ifx\csname#1@sh@#2@sel\endcsname\relax
    \bbl@afterelse\bbl@scndcs
  \else
    \bbl@afterfi\csname#1@sh@#2@sel\endcsname
  \fi}
%    \end{macrocode}
%
% \macro{\active@prefix}
% Used in the expansion of active
% characters has a function similar to |\OT1-cmd| in that it
% |\protect|s the active character whenever |\protect| is \emph{not}
% |\@typeset@protect|.
%
%    \begin{macrocode}
\begingroup
\bbl@ifunset{ifincsname}
  {\gdef\active@prefix#1{%
     \ifx\protect\@typeset@protect
     \else
       \ifx\protect\@unexpandable@protect
         \noexpand#1%
       \else
         \protect#1%
       \fi
       \expandafter\@gobble
     \fi}}
  {\gdef\active@prefix#1{%
     \ifincsname
       \string#1%
       \expandafter\@gobble
     \else
       \ifx\protect\@typeset@protect
       \else
         \ifx\protect\@unexpandable@protect
           \noexpand#1%
         \else
           \protect#1%
         \fi
         \expandafter\expandafter\expandafter\@gobble
       \fi
     \fi}}
\endgroup
%    \end{macrocode}
%
% \macro{if@safe@actives}
%
% In some circumstances it is necessary to be able to reset the
% shorthand to its ‘normal’ value (usually the character with catcode
% ‘other’) on the fly.
%
%    \begin{macrocode}
\newif\if@safe@actives
\@safe@activesfalse
%    \end{macrocode}
%
% \macro{\bbl@restore@actives}
% When the output routine kicks in while the active characters were
% made “safe” this must be undone in the headers to prevent
% unexpected typeset results. 
%
%    \begin{macrocode}
\def\bbl@restore@actives{\if@safe@actives\@safe@activesfalse\fi}
%    \end{macrocode}
%
% \macro{\bbl@activate}
% \macro{\bbl@deactivate}
%
% Used to change the definition of an active character to
% expand to |\active@char<char>| in the case of |\bbl@activate|,
% or |\normal@char<char>| in the case of |\bbl@deactivate|.
%
%    \begin{macrocode}
\chardef\bbl@activated\z@
\def\bbl@activate#1{%
  \chardef\bbl@activated\@ne
  \bbl@withactive{\expandafter\let\expandafter}#1%
    \csname bbl@active@\string#1\endcsname}
\def\bbl@deactivate#1{%
  \chardef\bbl@activated\tw@
  \bbl@withactive{\expandafter\let\expandafter}#1%
    \csname bbl@normal@\string#1\endcsname}
%    \end{macrocode}
%
% \macro{\bbl@firstcs}
% \macro{\bbl@scndcs}
%
% These macros are used only as a trick when declaring shorthands.
%
%    \begin{macrocode}
\def\bbl@firstcs#1#2{\csname#1\endcsname}
\def\bbl@scndcs#1#2{\csname#2\endcsname}
%    \end{macrocode}
%
% \macro{\declare@shorthand}
% Used to declare a shorthand on a certain level.
%
%    \begin{macrocode}
\def\babel@texpdf#1#2#3#4{%
  \ifx\texorpdfstring\@undefined
    \textormath{#1}{#3}%
  \else
    \texorpdfstring{\textormath{#1}{#3}}{#2}%
    % \texorpdfstring{\textormath{#1}{#3}}{\textormath{#2}{#4}}%
  \fi}
%
\def\declare@shorthand#1#2{\@decl@short{#1}#2\@nil}
\def\@decl@short#1#2#3\@nil#4{%
  \def\bbl@tempa{#3}%
  \ifx\bbl@tempa\@empty
    \expandafter\let\csname #1@sh@\string#2@sel\endcsname\bbl@scndcs
    \bbl@ifunset{#1@sh@\string#2@}{}%
      {\def\bbl@tempa{#4}%
       \expandafter\ifx\csname#1@sh@\string#2@\endcsname\bbl@tempa
       \else
         \bbl@info
           {Redefining #1 shorthand \string#2\\%
            in language \CurrentOption}%
       \fi}%
    \@namedef{#1@sh@\string#2@}{#4}%
  \else
    \expandafter\let\csname #1@sh@\string#2@sel\endcsname\bbl@firstcs
    \bbl@ifunset{#1@sh@\string#2@\string#3@}{}%
      {\def\bbl@tempa{#4}%
       \expandafter\ifx\csname#1@sh@\string#2@\string#3@\endcsname\bbl@tempa
       \else
         \bbl@info
           {Redefining #1 shorthand \string#2\string#3\\%
            in language \CurrentOption}%
       \fi}%
    \@namedef{#1@sh@\string#2@\string#3@}{#4}%
  \fi}
%    \end{macrocode}
%
% \macro{\textormath}
% Some of the shorthands that will be declared by the language
% definition files have to be usable in both text and mathmode. 
%
%    \begin{macrocode}
\def\textormath{%
  \ifmmode
    \expandafter\@secondoftwo
  \else
    \expandafter\@firstoftwo
  \fi}
%    \end{macrocode}
%
% \macro{\user@group}
% \macro{\language@group}
% \macro{\system@group}
% The current concept of `shorthands' supports three levels or groups
% of shorthands. 
%
%    \begin{macrocode}
\def\user@group{user}
\def\language@group{english}
\def\system@group{system}
%    \end{macrocode}
%
% \macro{\useshorthands}
%
% This is the user level macro. 
%
%    \begin{macrocode}
\def\useshorthands{%
  \@ifstar\bbl@usesh@s{\bbl@usesh@x{}}}
\def\bbl@usesh@s#1{%
  \bbl@usesh@x
    {\AddBabelHook{babel-sh-\string#1}{afterextras}{\bbl@activate{#1}}}%
    {#1}}
\def\bbl@usesh@x#1#2{%
  \bbl@ifshorthand{#2}%
    {\def\user@group{user}%
     \initiate@active@char{#2}%
     #1%
     \bbl@activate{#2}}%
    {\bbl@error{shorthand-is-off}{}{#2}{}}}
%    \end{macrocode}
%
% \macro{\defineshorthand}
%
% Currently we only support two groups of user level shorthands,
% named internally |user| and |user@<language>| (language-dependent
% user shorthands). 
%
%    \begin{macrocode}
\def\user@language@group{user@\language@group}
\def\bbl@set@user@generic#1#2{%
  \bbl@ifunset{user@generic@active#1}%
    {\bbl@active@def#1\user@language@group{user@active}{user@generic@active}%
     \bbl@active@def#1\user@group{user@generic@active}{language@active}%
     \expandafter\edef\csname#2@sh@#1@@\endcsname{%
       \expandafter\noexpand\csname normal@char#1\endcsname}%
     \expandafter\edef\csname#2@sh@#1@\string\protect@\endcsname{%
       \expandafter\noexpand\csname user@active#1\endcsname}}%
  \@empty}
\newcommand\defineshorthand[3][user]{%
  \edef\bbl@tempa{\zap@space#1 \@empty}%
  \bbl@for\bbl@tempb\bbl@tempa{%
    \if*\expandafter\@car\bbl@tempb\@nil
      \edef\bbl@tempb{user@\expandafter\@gobble\bbl@tempb}%
      \@expandtwoargs
        \bbl@set@user@generic{\expandafter\string\@car#2\@nil}\bbl@tempb
    \fi
    \declare@shorthand{\bbl@tempb}{#2}{#3}}}
%    \end{macrocode}
%
% \macro{\languageshorthands}
%
% A user level command to change the language from which shorthands are
% used. 
%
%    \begin{macrocode}
\def\languageshorthands#1{%
  \bbl@ifsamestring{none}{#1}{}{%
    \bbl@once{short-\localename-#1}{%
      \bbl@info{'\localename' activates '#1' shorthands.\\Reported}}}%
  \def\language@group{#1}}
%    \end{macrocode}
%
% \macro{\aliasshorthand}
%
% \textit{Deprecated}. 
%
%    \begin{macrocode}
\def\aliasshorthand#1#2{%
  \bbl@ifshorthand{#2}%
    {\expandafter\ifx\csname active@char\string#2\endcsname\relax
       \ifx\document\@notprerr
         \@notshorthand{#2}%
       \else
         \initiate@active@char{#2}%
         \bbl@ccarg\let{active@char\string#2}{active@char\string#1}%
         \bbl@ccarg\let{normal@char\string#2}{normal@char\string#1}%
         \bbl@activate{#2}%
       \fi
     \fi}%
    {\bbl@error{shorthand-is-off}{}{#2}{}}}
%    \end{macrocode}
%
% \macro{\@notshorthand}
%
%    \begin{macrocode}
\def\@notshorthand#1{\bbl@error{not-a-shorthand}{#1}{}{}}
%    \end{macrocode}
%
% \macro{\shorthandon}
% \macro{\shorthandoff}
%
% The first level definition of these macros just passes the
% argument on to |\bbl@switch@sh|, adding |\@nil| at the end to
% denote the end of the list of characters.
%
%    \begin{macrocode}
\newcommand*\shorthandon[1]{\bbl@switch@sh\@ne#1\@nnil}
\DeclareRobustCommand*\shorthandoff{%
  \@ifstar{\bbl@shorthandoff\tw@}{\bbl@shorthandoff\z@}}
\def\bbl@shorthandoff#1#2{\bbl@switch@sh#1#2\@nnil}
%    \end{macrocode}
%
% \macro{\bbl@switch@sh}
%
% The macro |\bbl@switch@sh| takes the list of characters apart one
% by one and subsequently switches the category code of the
% shorthand character according to the first argument of
% |\bbl@switch@sh|.
%
%    \begin{macrocode}
\def\bbl@switch@sh#1#2{%
  \ifx#2\@nnil\else
    \bbl@ifunset{bbl@active@\string#2}%
      {\bbl@error{not-a-shorthand-b}{}{#2}{}}%
      {\ifcase#1%   off, on, off*
         \catcode`#212\relax
       \or
         \catcode`#2\active
         \bbl@ifunset{bbl@shdef@\string#2}%
           {}%
           {\bbl@withactive{\expandafter\let\expandafter}#2%
              \csname bbl@shdef@\string#2\endcsname
            \bbl@csarg\let{shdef@\string#2}\relax}%
         \ifcase\bbl@activated\or
           \bbl@activate{#2}%
         \else
           \bbl@deactivate{#2}%
         \fi
       \or
         \bbl@ifunset{bbl@shdef@\string#2}%
           {\bbl@withactive{\bbl@csarg\let{shdef@\string#2}}#2}%
           {}%
         \csname bbl@oricat@\string#2\endcsname
         \csname bbl@oridef@\string#2\endcsname
       \fi}%
    \bbl@afterfi\bbl@switch@sh#1%
  \fi}
%    \end{macrocode}
%
% Note the value is that at the expansion time; e.g., in the preamble
% shorthands are usually deactivated.
%
%    \begin{macrocode}
\def\babelshorthand{\active@prefix\babelshorthand\bbl@putsh}
\def\bbl@putsh#1{%
  \bbl@ifunset{bbl@active@\string#1}%
     {\bbl@putsh@i#1\@empty\@nnil}%
     {\csname bbl@active@\string#1\endcsname}}
\def\bbl@putsh@i#1#2\@nnil{%
  \csname\language@group @sh@\string#1@%
    \ifx\@empty#2\else\string#2@\fi\endcsname}
%    \end{macrocode}
%
% You may want to test if a character is a shorthand. Note it does not
% test whether the shorthand is on or off.
%
%    \begin{macrocode}
\newcommand\ifbabelshorthand[3]{\bbl@ifunset{bbl@active@\string#1}{#3}{#2}}
%    \end{macrocode}
%
% \macro{\bbl@prim@s}
% \macro{\bbl@pr@m@s}
%
% One of the internal macros that are involved in substituting |\prime|
% for each right quote in mathmode is |\prim@s|. 
%
%    \begin{macrocode}
\def\bbl@prim@s{%
  \prime\futurelet\@let@token\bbl@pr@m@s}
\def\bbl@if@primes#1#2{%
  \ifx#1\@let@token
    \expandafter\@firstoftwo
  \else\ifx#2\@let@token
    \bbl@afterelse\expandafter\@firstoftwo
  \else
    \bbl@afterfi\expandafter\@secondoftwo
  \fi\fi}
\begingroup
  \catcode`\^=7  \catcode`\*=\active  \lccode`\*=`\^
  \catcode`\'=12 \catcode`\"=\active  \lccode`\"=`\'
  \lowercase{%
    \gdef\bbl@pr@m@s{%
      \bbl@if@primes"'%
        \pr@@@s
        {\bbl@if@primes*^\pr@@@t\egroup}}}
\endgroup
%    \end{macrocode}
%
% Usually the |~| is active and expands to \verb*=\penalty\@M\ =.
% When it is written to the \file{aux} file it is written
% expanded. To prevent that and to be able to use the character |~|
% as a start character for a shorthand, it is redefined here as a
% one character shorthand on system level. 
%
%    \begin{macrocode}
\initiate@active@char{~}
\declare@shorthand{system}{~}{\leavevmode\nobreak\ }
\bbl@activate{~}
%    \end{macrocode}
%
% \macro{\OT1dqpos}
% \macro{\T1dqpos}
% The position of the double quote character is different for the OT1
% and T1 encodings.
%
%    \begin{macrocode}
\expandafter\def\csname OT1dqpos\endcsname{127}
\expandafter\def\csname T1dqpos\endcsname{4}
%    \end{macrocode}
%
% When the macro |\f@encoding| is undefined (as it is in plain \TeX) we
% define it here to expand to \texttt{OT1}
%
%    \begin{macrocode}
\ifx\f@encoding\@undefined
  \def\f@encoding{OT1}
\fi
%    \end{macrocode}
%
% \subsection{Language attributes}
%
% Language attributes provide a means to give the user control over
% which features of the language definition files he wants to enable.
%
% \macro{\languageattribute}
%
% The macro |\languageattribute| checks whether its arguments are valid
% and then activates the selected language attribute.
%
%    \begin{macrocode}
\bbl@trace{Language attributes}
\newcommand\languageattribute[2]{%
  \def\bbl@tempc{#1}%
  \bbl@fixname\bbl@tempc
  \bbl@iflanguage\bbl@tempc{%
    \bbl@vforeach{#2}{%
      \ifx\bbl@known@attribs\@undefined
        \in@false
      \else
        \bbl@xin@{,\bbl@tempc-##1,}{,\bbl@known@attribs,}%
      \fi
      \ifin@
        \bbl@warning{%
          You have more than once selected the attribute '##1'\\%
          for language #1. Reported}%
      \else
        \bbl@info{Activated '##1' attribute for\\%
          '\bbl@tempc'. Reported}%
        \bbl@exp{%
          \\\bbl@add@list\\\bbl@known@attribs{\bbl@tempc-##1}}%
        \edef\bbl@tempa{\bbl@tempc-##1}%
        \expandafter\bbl@ifknown@ttrib\expandafter{\bbl@tempa}\bbl@attributes
          {\csname\bbl@tempc @attr@##1\endcsname}%
          {\bbl@error{unknown-attribute}{\bbl@tempc}{##1}{}}%
     \fi}}}
\@onlypreamble\languageattribute
%    \end{macrocode}
%
% \macro{\bbl@declare@ttribute}
% This command adds the new language/attribute combination to the list
% of known attributes.
%
%    \begin{macrocode}
\def\bbl@declare@ttribute#1#2#3{%
  \bbl@xin@{,#2,}{,\BabelModifiers,}%
  \ifin@
    \AfterBabelLanguage{#1}{\languageattribute{#1}{#2}}%
  \fi
  \bbl@add@list\bbl@attributes{#1-#2}%
  \expandafter\def\csname#1@attr@#2\endcsname{#3}}
%    \end{macrocode}
%
% \macro{\bbl@ifattributeset}
%
% This internal macro has 4 arguments. It can be used to interpret
% \TeX\ code based on whether a certain attribute was set. T
%
%    \begin{macrocode}
\def\bbl@ifattributeset#1#2#3#4{%
  \ifx\bbl@known@attribs\@undefined
    \in@false
  \else
    \bbl@xin@{,#1-#2,}{,\bbl@known@attribs,}%
  \fi
  \ifin@
    \bbl@afterelse#3%
  \else
    \bbl@afterfi#4%
  \fi}
%    \end{macrocode}
%
% \macro{\bbl@ifknown@ttrib}
% An internal macro to check whether a given language/attribute is
% known. 
%
%    \begin{macrocode}
\def\bbl@ifknown@ttrib#1#2{%
  \let\bbl@tempa\@secondoftwo
  \bbl@loopx\bbl@tempb{#2}{%
    \expandafter\in@\expandafter{\expandafter,\bbl@tempb,}{,#1,}%
    \ifin@
      \let\bbl@tempa\@firstoftwo
    \else
    \fi}%
  \bbl@tempa}
%    \end{macrocode}
%
% \macro{\bbl@clear@ttribs}
% This macro removes all the attribute code from \LaTeX’s memory at
% |\begin{document}| time (if any is present).
%
%    \begin{macrocode}
\def\bbl@clear@ttribs{%
  \ifx\bbl@attributes\@undefined\else
    \bbl@loopx\bbl@tempa{\bbl@attributes}{%
      \expandafter\bbl@clear@ttrib\bbl@tempa.}%
    \let\bbl@attributes\@undefined
  \fi}
\def\bbl@clear@ttrib#1-#2.{%
  \expandafter\let\csname#1@attr@#2\endcsname\@undefined}
\AtBeginDocument{\bbl@clear@ttribs}
%    \end{macrocode}
%
% \subsection{Support for saving and redefining macros}
%
% To save the meaning of control sequences using |\babel@save|, we use
% temporary control sequences. 
%
% \macro{\babel@savecnt}
% \macro{\babel@beginsave}
% The initialization of a new save cycle: reset the counter to zero.
%
%    \begin{macrocode}
\bbl@trace{Macros for saving definitions}
\def\babel@beginsave{\babel@savecnt\z@}
%    \end{macrocode}
%
% Before it’s forgotten, allocate the counter and initialize all.
%
%    \begin{macrocode}
\newcount\babel@savecnt
\babel@beginsave
%    \end{macrocode}
%
% \macro{\babel@save}
% \macro{\babel@savevariable}
%
% The macro |\babel@save|\meta{csname} saves the current meaning of the
% control sequence \meta{csname} to |\originalTeX| (which has to be
% expandable, i.e., you shouldn't let it to \cs{relax}). 
%
%    \begin{macrocode}
\def\babel@save#1{%
  \def\bbl@tempa{{,#1,}}% Clumsy, for Plain
  \expandafter\bbl@add\expandafter\bbl@tempa\expandafter{%
    \expandafter{\expandafter,\bbl@savedextras,}}%
  \expandafter\in@\bbl@tempa
  \ifin@\else
    \bbl@add\bbl@savedextras{,#1,}%
    \bbl@carg\let{babel@\number\babel@savecnt}#1\relax
    \toks@\expandafter{\originalTeX\let#1=}%
    \bbl@exp{%
      \def\\\originalTeX{\the\toks@\<babel@\number\babel@savecnt>\relax}}%
    \advance\babel@savecnt\@ne
  \fi}
\def\babel@savevariable#1{%
  \toks@\expandafter{\originalTeX #1=}%
  \bbl@exp{\def\\\originalTeX{\the\toks@\the#1\relax}}}
%    \end{macrocode}
%
% \macro{\bbl@redefine}
%
% To redefine a command, we save the old meaning of the macro. 
%
%    \begin{macrocode}
\def\bbl@redefine#1{%
  \edef\bbl@tempa{\bbl@stripslash#1}%
  \expandafter\let\csname org@\bbl@tempa\endcsname#1%
  \expandafter\def\csname\bbl@tempa\endcsname}
\@onlypreamble\bbl@redefine
%    \end{macrocode}
%
% \macro{\bbl@redefine@long}
% This version of |\babel@redefine| can be used to redefine |\long|
% commands such as |\ifthenelse|.
%
%    \begin{macrocode}
\def\bbl@redefine@long#1{%
  \edef\bbl@tempa{\bbl@stripslash#1}%
  \expandafter\let\csname org@\bbl@tempa\endcsname#1%
  \long\expandafter\def\csname\bbl@tempa\endcsname}
\@onlypreamble\bbl@redefine@long
%    \end{macrocode}
%
% \macro{\bbl@redefinerobust}
% For commands that are redefined, but which \textit{might} be
% robust we need a slightly more intelligent macro. 
%
%    \begin{macrocode}
\def\bbl@redefinerobust#1{%
  \edef\bbl@tempa{\bbl@stripslash#1}%
  \bbl@ifunset{\bbl@tempa\space}%
    {\expandafter\let\csname org@\bbl@tempa\endcsname#1%
     \bbl@exp{\def\\#1{\\\protect\<\bbl@tempa\space>}}}%
    {\bbl@exp{\let\<org@\bbl@tempa>\<\bbl@tempa\space>}}%
    \@namedef{\bbl@tempa\space}}
\@onlypreamble\bbl@redefinerobust
%    \end{macrocode}
%
% \subsection{French spacing}
%
% \macro{\bbl@frenchspacing}
% \macro{\bbl@nonfrenchspacing}
% Some languages need to have |\frenchspacing| in effect. 
%
%    \begin{macrocode}
\def\bbl@frenchspacing{%
  \ifnum\the\sfcode`\.=\@m
    \let\bbl@nonfrenchspacing\relax
  \else
    \frenchspacing
    \let\bbl@nonfrenchspacing\nonfrenchspacing
  \fi}
\let\bbl@nonfrenchspacing\nonfrenchspacing
%    \end{macrocode}
%
% \subsection{Hyphens}
%
% \macro{\babelhyphenation}
%
% This macro saves hyphenation exceptions. 
%
%    \begin{macrocode}
\bbl@trace{Hyphens}
\@onlypreamble\babelhyphenation
\AtEndOfPackage{%
  \newcommand\babelhyphenation[2][\@empty]{%
    \ifx\bbl@hyphenation@\relax
      \let\bbl@hyphenation@\@empty
    \fi
    \ifx\bbl@hyphlist\@empty\else
      \bbl@warning{%
        You must not intermingle \string\selectlanguage\space and\\%
        \string\babelhyphenation\space or some exceptions will not\\%
        be taken into account. Reported}%
    \fi
    \ifx\@empty#1%
      \protected@edef\bbl@hyphenation@{\bbl@hyphenation@\space#2}%
    \else
      \bbl@vforeach{#1}{%
        \def\bbl@tempa{##1}%
        \bbl@fixname\bbl@tempa
        \bbl@iflanguage\bbl@tempa{%
          \bbl@csarg\protected@edef{hyphenation@\bbl@tempa}{%
            \bbl@ifunset{bbl@hyphenation@\bbl@tempa}%
              {}%
              {\csname bbl@hyphenation@\bbl@tempa\endcsname\space}%
            #2}}}%
    \fi}}
%    \end{macrocode}
%
% \macro{\bbl@allowhyphens}
%
% This macro makes hyphenation possible. 
%
%    \begin{macrocode}
\def\bbl@allowhyphens{\ifvmode\else\nobreak\hskip\z@skip\fi}
\def\bbl@t@one{T1}
\def\allowhyphens{\ifx\cf@encoding\bbl@t@one\else\bbl@allowhyphens\fi}
%    \end{macrocode}
%
% \macro{\babelhyphen}
%
% Macros to insert common hyphens.
%
%    \begin{macrocode}
\newcommand\babelnullhyphen{\char\hyphenchar\font}
\def\babelhyphen{\active@prefix\babelhyphen\bbl@hyphen}
\def\bbl@hyphen{%
  \@ifstar{\bbl@hyphen@i @}{\bbl@hyphen@i\@empty}}
\def\bbl@hyphen@i#1#2{%
  \lowercase{\bbl@ifunset{bbl@hy@#1#2\@empty}}%
    {\csname bbl@#1usehyphen\endcsname{\discretionary{#2}{}{#2}}}%
    {\lowercase{\csname bbl@hy@#1#2\@empty\endcsname}}}
%    \end{macrocode}
%
% The following two commands are used to wrap the “hyphen” and
% set the behavior of the rest of the word.
%
%    \begin{macrocode}
\def\bbl@usehyphen#1{%
  \leavevmode
  \ifdim\lastskip>\z@\mbox{#1}\else\nobreak#1\fi
  \nobreak\hskip\z@skip}
\def\bbl@@usehyphen#1{%
  \leavevmode\ifdim\lastskip>\z@\mbox{#1}\else#1\fi}
%    \end{macrocode}
%
% The following macro inserts the hyphen char.
%
%    \begin{macrocode}
\def\bbl@hyphenchar{%
  \ifnum\hyphenchar\font=\m@ne
    \babelnullhyphen
  \else
    \char\hyphenchar\font
  \fi}
%    \end{macrocode}
%
% Finally, we define the hyphen “types”. Their names will not change,
% so you may use them in |ldf|’s. After a space, the |\mbox| in
% |\bbl@hy@nobreak| is redundant.
%
%    \begin{macrocode}
\def\bbl@hy@soft{\bbl@usehyphen{\discretionary{\bbl@hyphenchar}{}{}}}
\def\bbl@hy@@soft{\bbl@@usehyphen{\discretionary{\bbl@hyphenchar}{}{}}}
\def\bbl@hy@hard{\bbl@usehyphen\bbl@hyphenchar}
\def\bbl@hy@@hard{\bbl@@usehyphen\bbl@hyphenchar}
\def\bbl@hy@nobreak{\bbl@usehyphen{\mbox{\bbl@hyphenchar}}}
\def\bbl@hy@@nobreak{\mbox{\bbl@hyphenchar}}
\def\bbl@hy@repeat{%
  \bbl@usehyphen{%
    \discretionary{\bbl@hyphenchar}{\bbl@hyphenchar}{\bbl@hyphenchar}}}
\def\bbl@hy@@repeat{%
  \bbl@@usehyphen{%
    \discretionary{\bbl@hyphenchar}{\bbl@hyphenchar}{\bbl@hyphenchar}}}
\def\bbl@hy@empty{\hskip\z@skip}
\def\bbl@hy@@empty{\discretionary{}{}{}}
%    \end{macrocode}
%
% \macro{\bbl@disc}
% For some languages the macro |\bbl@disc| is used to ease the
% insertion of discretionaries for letters that behave `abnormally' at
% a breakpoint.
%
%    \begin{macrocode}
\def\bbl@disc#1#2{\nobreak\discretionary{#2-}{}{#1}\bbl@allowhyphens}
%    \end{macrocode}
%
% \subsection{Multiencoding strings}
%
% The aim following commands is to provide a common interface for
% strings in several encodings. 
%
%  \paragraph{Tools}
%
% But first, a tool. It makes global a local variable. This is not the
% best solution, but it works.
%
%    \begin{macrocode}
\bbl@trace{Multiencoding strings}
\def\bbl@toglobal#1{\global\let#1#1}
%    \end{macrocode}
%
% \paragraph{Main command} This is the main command. 
%
%    \begin{macrocode}
\@onlypreamble\StartBabelCommands
\def\StartBabelCommands{%
  \begingroup
  \@tempcnta="7F
  \def\bbl@tempa{%
    \ifnum\@tempcnta>"FF\else
      \catcode\@tempcnta=11
      \advance\@tempcnta\@ne
      \expandafter\bbl@tempa
    \fi}%
  \bbl@tempa
  <@Macros local to Plain BabelCommands@>
  \def\bbl@provstring##1##2{%
    \providecommand##1{##2}%
    \bbl@toglobal##1}%
  \global\let\bbl@scafter\@empty
  \let\StartBabelCommands\bbl@startcmds
  \ifx\BabelLanguages\relax
     \let\BabelLanguages\CurrentOption
  \fi
  \begingroup
  \let\bbl@screset\@nnil % local flag - disable 1st stopcommands
  \StartBabelCommands}
\def\bbl@startcmds{%
  \ifx\bbl@screset\@nnil\else
    \bbl@usehooks{stopcommands}{}%
  \fi
  \endgroup
  \begingroup
  \@ifstar
    {\ifx\bbl@opt@strings\@nnil
       \let\bbl@opt@strings\BabelStringsDefault
     \fi
     \bbl@startcmds@i}%
    \bbl@startcmds@i}
\def\bbl@startcmds@i#1#2{%
  \edef\bbl@L{\zap@space#1 \@empty}%
  \edef\bbl@G{\zap@space#2 \@empty}%
  \bbl@startcmds@ii}
\let\bbl@startcommands\StartBabelCommands
%    \end{macrocode}
%
% Parse the encoding info to get the label, input, and font parts.
%
%    \begin{macrocode}
\newcommand\bbl@startcmds@ii[1][\@empty]{%
  \let\SetString\@gobbletwo
  \let\bbl@stringdef\@gobbletwo
  \let\AfterBabelCommands\@gobble
  \ifx\@empty#1%
    \def\bbl@sc@label{generic}%
    \def\bbl@encstring##1##2{%
      \ProvideTextCommandDefault##1{##2}%
      \bbl@toglobal##1%
      \expandafter\bbl@toglobal\csname\string?\string##1\endcsname}%
    \let\bbl@sctest\in@true
  \else
    \let\bbl@sc@charset\space % <- zapped below
    \let\bbl@sc@fontenc\space % <-   "      "
    \def\bbl@tempa##1=##2\@nil{%
      \bbl@csarg\edef{sc@\zap@space##1 \@empty}{##2 }}%
    \bbl@vforeach{label=#1}{\bbl@tempa##1\@nil}%
    \def\bbl@tempa##1 ##2{% space -> comma
      ##1%
      \ifx\@empty##2\else\ifx,##1,\else,\fi\bbl@afterfi\bbl@tempa##2\fi}%
    \edef\bbl@sc@fontenc{\expandafter\bbl@tempa\bbl@sc@fontenc\@empty}%
    \edef\bbl@sc@label{\expandafter\zap@space\bbl@sc@label\@empty}%
    \edef\bbl@sc@charset{\expandafter\zap@space\bbl@sc@charset\@empty}%
    \def\bbl@encstring##1##2{%
      \bbl@foreach\bbl@sc@fontenc{%
        \bbl@ifunset{T@####1}%
          {}%
          {\ProvideTextCommand##1{####1}{##2}%
           \bbl@toglobal##1%
           \expandafter
           \bbl@toglobal\csname####1\string##1\endcsname}}}%
    \def\bbl@sctest{%
      \bbl@xin@{,\bbl@opt@strings,}{,\bbl@sc@label,\bbl@sc@fontenc,}}%
  \fi
  \ifx\bbl@opt@strings\@nnil         % i.e., no strings key -> defaults
  \else\ifx\bbl@opt@strings\relax    % i.e., strings=encoded
    \let\AfterBabelCommands\bbl@aftercmds
    \let\SetString\bbl@setstring
    \let\bbl@stringdef\bbl@encstring
  \else       % i.e., strings=value
  \bbl@sctest
  \ifin@
    \let\AfterBabelCommands\bbl@aftercmds
    \let\SetString\bbl@setstring
    \let\bbl@stringdef\bbl@provstring
  \fi\fi\fi
  \bbl@scswitch
  \ifx\bbl@G\@empty
    \def\SetString##1##2{%
      \bbl@error{missing-group}{##1}{}{}}%
  \fi
  \ifx\@empty#1%
    \bbl@usehooks{defaultcommands}{}%
  \else
    \@expandtwoargs
    \bbl@usehooks{encodedcommands}{{\bbl@sc@charset}{\bbl@sc@fontenc}}%
  \fi}
%    \end{macrocode}
%
% There are two versions of |\bbl@scswitch|. 
%
%   \begin{macrocode}
\def\bbl@forlang#1#2{%
  \bbl@for#1\bbl@L{%
    \bbl@xin@{,#1,}{,\BabelLanguages,}%
    \ifin@#2\relax\fi}}
\def\bbl@scswitch{%
  \bbl@forlang\bbl@tempa{%
    \ifx\bbl@G\@empty\else
      \ifx\SetString\@gobbletwo\else
        \edef\bbl@GL{\bbl@G\bbl@tempa}%
        \bbl@xin@{,\bbl@GL,}{,\bbl@screset,}%
        \ifin@\else
          \global\expandafter\let\csname\bbl@GL\endcsname\@undefined
          \xdef\bbl@screset{\bbl@screset,\bbl@GL}%
        \fi
      \fi
    \fi}}
\AtEndOfPackage{%
  \def\bbl@forlang#1#2{\bbl@for#1\bbl@L{\bbl@ifunset{date#1}{}{#2}}}%
  \let\bbl@scswitch\relax}
\@onlypreamble\EndBabelCommands
\def\EndBabelCommands{%
  \bbl@usehooks{stopcommands}{}%
  \endgroup
  \endgroup
  \bbl@scafter}
\let\bbl@endcommands\EndBabelCommands
%    \end{macrocode}
%
% Now we define commands to be used inside |\StartBabelCommands|.
%
% \paragraph{Strings} The following macro is the actual definition
% of |\SetString| when it is “active”,
%
%    \begin{macrocode}
\def\bbl@setstring#1#2{% e.g., \prefacename{<string>}
  \bbl@forlang\bbl@tempa{%
    \edef\bbl@LC{\bbl@tempa\bbl@stripslash#1}%
    \bbl@ifunset{\bbl@LC}% e.g., \germanchaptername
      {\bbl@exp{%
         \global\\\bbl@add\<\bbl@G\bbl@tempa>{\\\bbl@scset\\#1\<\bbl@LC>}}}%
      {}%
    \def\BabelString{#2}%
    \bbl@usehooks{stringprocess}{}%
    \expandafter\bbl@stringdef
      \csname\bbl@LC\expandafter\endcsname\expandafter{\BabelString}}}
%    \end{macrocode}
%
% A little auxiliary command sets the string. Formerly used with
% casing. Very likely no longer necessary, although it’s used in
% |\setlocalecaption|.
%
%    \begin{macrocode}
\def\bbl@scset#1#2{\def#1{#2}}
%    \end{macrocode}
%
% Define |\SetStringLoop|, which is actually set inside
% |\StartBabelCommands|.
%
%    \begin{macrocode}
%<<*Macros local to Plain BabelCommands>>
\def\SetStringLoop##1##2{%
    \def\bbl@templ####1{\expandafter\noexpand\csname##1\endcsname}%
    \count@\z@
    \bbl@loop\bbl@tempa{##2}{% empty items and spaces are ok
      \advance\count@\@ne
      \toks@\expandafter{\bbl@tempa}%
      \bbl@exp{%
        \\\SetString\bbl@templ{\romannumeral\count@}{\the\toks@}%
        \count@=\the\count@\relax}}}%
%<</Macros local to Plain BabelCommands>>
%    \end{macrocode}
%
% \paragraph{Delaying code} Now the definition of |\AfterBabelCommands|
% when it is activated.
%
%    \begin{macrocode}
\def\bbl@aftercmds#1{%
  \toks@\expandafter{\bbl@scafter#1}%
  \xdef\bbl@scafter{\the\toks@}}
%    \end{macrocode}
%
% \paragraph{Case mapping}
%
% The command |\SetCase| is deprecated (at least in \LaTeX). This is
% the original definition.
%
%    \begin{macrocode}
\def\bbl@patchuclc{%
  \global\let\bbl@patchuclc\relax
  \g@addto@macro\@uclclist{\reserved@b{\reserved@b\bbl@uclc}}%
  \gdef\bbl@uclc##1{%
    \let\bbl@encoded\bbl@encoded@uclc
    \bbl@ifunset{\languagename @bbl@uclc}% and resumes it
      {##1}%
      {\let\bbl@tempa##1\relax % Used by LANG@bbl@uclc
       \csname\languagename @bbl@uclc\endcsname}%
    {\bbl@tolower\@empty}{\bbl@toupper\@empty}}%
  \gdef\bbl@tolower{\csname\languagename @bbl@lc\endcsname}%
  \gdef\bbl@toupper{\csname\languagename @bbl@uc\endcsname}}
%<<*Macros local to Plain BabelCommands>>
  \newcommand\SetCase[3][]{%
    \bbl@patchuclc
    \bbl@forlang\bbl@tempa{%
      \bbl@carg\bbl@encstring{\bbl@tempa @bbl@uclc}{\bbl@tempa##1}%
      \bbl@carg\bbl@encstring{\bbl@tempa @bbl@uc}{##2}%
      \bbl@carg\bbl@encstring{\bbl@tempa @bbl@lc}{##3}}}%
%<</Macros local to Plain BabelCommands>>
%    \end{macrocode}
%
% Macros to deal with case mapping for hyphenation. 
%
%    \begin{macrocode}
%<<*Macros local to Plain BabelCommands>>
  \newcommand\SetHyphenMap[1]{%
    \bbl@forlang\bbl@tempa{%
      \expandafter\bbl@stringdef
        \csname\bbl@tempa @bbl@hyphenmap\endcsname{##1}}}%
%<</Macros local to Plain BabelCommands>>
%    \end{macrocode}
%
% There are 3 helper macros which do most of the work for you.
%
%    \begin{macrocode}
\newcommand\BabelLower[2]{% one to one.
  \ifnum\lccode#1=#2\else
    \babel@savevariable{\lccode#1}%
    \lccode#1=#2\relax
  \fi}
\newcommand\BabelLowerMM[4]{% many-to-many
  \@tempcnta=#1\relax
  \@tempcntb=#4\relax
  \def\bbl@tempa{%
    \ifnum\@tempcnta>#2\else
      \@expandtwoargs\BabelLower{\the\@tempcnta}{\the\@tempcntb}%
      \advance\@tempcnta#3\relax
      \advance\@tempcntb#3\relax
      \expandafter\bbl@tempa
    \fi}%
  \bbl@tempa}
\newcommand\BabelLowerMO[4]{% many-to-one
  \@tempcnta=#1\relax
  \def\bbl@tempa{%
    \ifnum\@tempcnta>#2\else
      \@expandtwoargs\BabelLower{\the\@tempcnta}{#4}%
      \advance\@tempcnta#3
      \expandafter\bbl@tempa
    \fi}%
  \bbl@tempa}
%    \end{macrocode}
%
% Initial setup to provide a default behavior if |hyphenmap| is not
% set.
%
%    \begin{macrocode}
\AtEndOfPackage{%
  \ifx\bbl@opt@hyphenmap\@undefined
    \bbl@xin@{,}{\bbl@language@opts}%
    \chardef\bbl@opt@hyphenmap\ifin@4\else\@ne\fi
  \fi}
%    \end{macrocode}
%
% \subsection{Tailor captions}
%
% A general tool for resetting the caption names with a unique
% interface.
%
%    \begin{macrocode}
\newcommand\setlocalecaption{%
  \@ifstar\bbl@setcaption@s\bbl@setcaption@x}
\def\bbl@setcaption@x#1#2#3{%  language caption-name string
  \bbl@trim@def\bbl@tempa{#2}%
  \bbl@xin@{.template}{\bbl@tempa}%
  \ifin@
    \bbl@ini@captions@template{#3}{#1}%
  \else
    \edef\bbl@tempd{%
      \expandafter\expandafter\expandafter
      \strip@prefix\expandafter\meaning\csname captions#1\endcsname}%
    \bbl@xin@
      {\expandafter\string\csname #2name\endcsname}%
      {\bbl@tempd}%
    \ifin@ % Renew caption
      \bbl@xin@{\string\bbl@scset}{\bbl@tempd}%
      \ifin@
        \bbl@exp{%
          \\\bbl@ifsamestring{\bbl@tempa}{\languagename}%
            {\\\bbl@scset\<#2name>\<#1#2name>}%
            {}}%
      \else % Old way converts to new way
        \bbl@ifunset{#1#2name}%
          {\bbl@exp{%
            \\\bbl@add\<captions#1>{\def\<#2name>{\<#1#2name>}}%
            \\\bbl@ifsamestring{\bbl@tempa}{\languagename}%
              {\def\<#2name>{\<#1#2name>}}%
              {}}}%
          {}%
      \fi
    \else
      \bbl@xin@{\string\bbl@scset}{\bbl@tempd}% New
      \ifin@ % New way
        \bbl@exp{%
          \\\bbl@add\<captions#1>{\\\bbl@scset\<#2name>\<#1#2name>}%
          \\\bbl@ifsamestring{\bbl@tempa}{\languagename}%
            {\\\bbl@scset\<#2name>\<#1#2name>}%
            {}}%
      \else  % Old way, but defined in the new way
        \bbl@exp{%
          \\\bbl@add\<captions#1>{\def\<#2name>{\<#1#2name>}}%
          \\\bbl@ifsamestring{\bbl@tempa}{\languagename}%
            {\def\<#2name>{\<#1#2name>}}%
            {}}%
      \fi%
    \fi
    \@namedef{#1#2name}{#3}%
    \toks@\expandafter{\bbl@captionslist}%
    \bbl@exp{\\\in@{\<#2name>}{\the\toks@}}%
    \ifin@\else
      \bbl@exp{\\\bbl@add\\\bbl@captionslist{\<#2name>}}%
      \bbl@toglobal\bbl@captionslist
    \fi
  \fi}
%    \end{macrocode}
%
% \subsection{Making glyphs available}
%
% This section makes a number of glyphs available that either do not
% exist in the \texttt{OT1} encoding and have to be `faked', or
% that are not accessible through \file{T1enc.def}.
%
% \macro{\set@low@box}
%
% The following macro is used to lower quotes to the same level as
% the comma.  
%
%    \begin{macrocode}
\bbl@trace{Macros related to glyphs}
\def\set@low@box#1{\setbox\tw@\hbox{,}\setbox\z@\hbox{#1}%
    \dimen\z@\ht\z@ \advance\dimen\z@ -\ht\tw@%
    \setbox\z@\hbox{\lower\dimen\z@ \box\z@}\ht\z@\ht\tw@ \dp\z@\dp\tw@}
%    \end{macrocode}
%
% \macro{\save@sf@q}
% The macro |\save@sf@q| is used to save and reset the current space
% factor.
%
%    \begin{macrocode}
\def\save@sf@q#1{\leavevmode
  \begingroup
    \edef\@SF{\spacefactor\the\spacefactor}#1\@SF
  \endgroup}
%    \end{macrocode}
%
% \subsubsection{Quotation marks}
%
% \macro{\quotedblbase}
% In the \texttt{OT1} encoding it is not
% available, therefore we make it available by lowering the normal
% open quote character to the baseline.
%
%    \begin{macrocode}
\ProvideTextCommand{\quotedblbase}{OT1}{%
  \save@sf@q{\set@low@box{\textquotedblright\/}%
    \box\z@\kern-.04em\bbl@allowhyphens}}
%    \end{macrocode}
%
% Make sure that when an encoding other than \texttt{OT1} or
% \texttt{T1} is used this glyph can still be typeset.
%
%    \begin{macrocode}
\ProvideTextCommandDefault{\quotedblbase}{%
  \UseTextSymbol{OT1}{\quotedblbase}}
%    \end{macrocode}
%
% \macro{\quotesinglbase}
% We also need the single quote character at the baseline.
%
%    \begin{macrocode}
\ProvideTextCommand{\quotesinglbase}{OT1}{%
  \save@sf@q{\set@low@box{\textquoteright\/}%
    \box\z@\kern-.04em\bbl@allowhyphens}}
%    \end{macrocode}
%
% Make sure that when an encoding other than \texttt{OT1} or
% \texttt{T1} is used this glyph can still be typeset.
%
%    \begin{macrocode}
\ProvideTextCommandDefault{\quotesinglbase}{%
  \UseTextSymbol{OT1}{\quotesinglbase}}
%    \end{macrocode}
%
% \macro{\guillemetleft}
% \macro{\guillemetright}
% The guillemet characters are not available in \texttt{OT1}
% encoding.
%
%    \begin{macrocode}
\ProvideTextCommand{\guillemetleft}{OT1}{%
  \ifmmode
    \ll
  \else
    \save@sf@q{\nobreak
      \raise.2ex\hbox{$\scriptscriptstyle\ll$}\bbl@allowhyphens}%
  \fi}
\ProvideTextCommand{\guillemetright}{OT1}{%
  \ifmmode
    \gg
  \else
    \save@sf@q{\nobreak
      \raise.2ex\hbox{$\scriptscriptstyle\gg$}\bbl@allowhyphens}%
  \fi}
\ProvideTextCommand{\guillemotleft}{OT1}{%
  \ifmmode
    \ll
  \else
    \save@sf@q{\nobreak
      \raise.2ex\hbox{$\scriptscriptstyle\ll$}\bbl@allowhyphens}%
  \fi}
\ProvideTextCommand{\guillemotright}{OT1}{%
  \ifmmode
    \gg
  \else
    \save@sf@q{\nobreak
      \raise.2ex\hbox{$\scriptscriptstyle\gg$}\bbl@allowhyphens}%
  \fi}
%    \end{macrocode}
%
% Make sure that when an encoding other than \texttt{OT1} or
% \texttt{T1} is used these glyphs can still be typeset.
%
%    \begin{macrocode}
\ProvideTextCommandDefault{\guillemetleft}{%
  \UseTextSymbol{OT1}{\guillemetleft}}
\ProvideTextCommandDefault{\guillemetright}{%
  \UseTextSymbol{OT1}{\guillemetright}}
\ProvideTextCommandDefault{\guillemotleft}{%
  \UseTextSymbol{OT1}{\guillemotleft}}
\ProvideTextCommandDefault{\guillemotright}{%
  \UseTextSymbol{OT1}{\guillemotright}}
%    \end{macrocode}
%
% \macro{\guilsinglleft}
% \macro{\guilsinglright}
% The single guillemets are not available in \texttt{OT1} encoding.
% They are faked.
%
%    \begin{macrocode}
\ProvideTextCommand{\guilsinglleft}{OT1}{%
  \ifmmode
    <%
  \else
    \save@sf@q{\nobreak
      \raise.2ex\hbox{$\scriptscriptstyle<$}\bbl@allowhyphens}%
  \fi}
\ProvideTextCommand{\guilsinglright}{OT1}{%
  \ifmmode
    >%
  \else
    \save@sf@q{\nobreak
      \raise.2ex\hbox{$\scriptscriptstyle>$}\bbl@allowhyphens}%
  \fi}
%    \end{macrocode}
%
% Make sure that when an encoding other than \texttt{OT1} or
% \texttt{T1} is used these glyphs can still be typeset.
%
%    \begin{macrocode}
\ProvideTextCommandDefault{\guilsinglleft}{%
  \UseTextSymbol{OT1}{\guilsinglleft}}
\ProvideTextCommandDefault{\guilsinglright}{%
  \UseTextSymbol{OT1}{\guilsinglright}}
%    \end{macrocode}
%
% \subsubsection{Letters}
%
% \macro{\ij}
% \macro{\IJ}
% The dutch language uses the letter `ij'. 
%
%    \begin{macrocode}
\DeclareTextCommand{\ij}{OT1}{%
  i\kern-0.02em\bbl@allowhyphens j}
\DeclareTextCommand{\IJ}{OT1}{%
  I\kern-0.02em\bbl@allowhyphens J}
\DeclareTextCommand{\ij}{T1}{\char188}
\DeclareTextCommand{\IJ}{T1}{\char156}
%    \end{macrocode}
%
% Make sure that when an encoding other than \texttt{OT1} or
% \texttt{T1} is used these glyphs can still be typeset.
%
%    \begin{macrocode}
\ProvideTextCommandDefault{\ij}{%
  \UseTextSymbol{OT1}{\ij}}
\ProvideTextCommandDefault{\IJ}{%
  \UseTextSymbol{OT1}{\IJ}}
%    \end{macrocode}
%
% \macro{\dj}
% \macro{\DJ}
% The croatian language needs the letters |\dj| and |\DJ|,
%
%    \begin{macrocode}
\def\crrtic@{\hrule height0.1ex width0.3em}
\def\crttic@{\hrule height0.1ex width0.33em}
\def\ddj@{%
  \setbox0\hbox{d}\dimen@=\ht0
  \advance\dimen@1ex
  \dimen@.45\dimen@
  \dimen@ii\expandafter\rem@pt\the\fontdimen\@ne\font\dimen@
  \advance\dimen@ii.5ex
  \leavevmode\rlap{\raise\dimen@\hbox{\kern\dimen@ii\vbox{\crrtic@}}}}
\def\DDJ@{%
  \setbox0\hbox{D}\dimen@=.55\ht0
  \dimen@ii\expandafter\rem@pt\the\fontdimen\@ne\font\dimen@
  \advance\dimen@ii.15ex %            correction for the dash position
  \advance\dimen@ii-.15\fontdimen7\font %     correction for cmtt font
  \dimen\thr@@\expandafter\rem@pt\the\fontdimen7\font\dimen@
  \leavevmode\rlap{\raise\dimen@\hbox{\kern\dimen@ii\vbox{\crttic@}}}}
%
\DeclareTextCommand{\dj}{OT1}{\ddj@ d}
\DeclareTextCommand{\DJ}{OT1}{\DDJ@ D}
%    \end{macrocode}
%
% Make sure that when an encoding other than \texttt{OT1} or
% \texttt{T1} is used these glyphs can still be typeset.
%
%    \begin{macrocode}
\ProvideTextCommandDefault{\dj}{%
  \UseTextSymbol{OT1}{\dj}}
\ProvideTextCommandDefault{\DJ}{%
  \UseTextSymbol{OT1}{\DJ}}
%    \end{macrocode}
%
% \macro{\SS}
% For the \texttt{T1} encoding |\SS| is defined and selects a specific
% glyph from the font, but for other encodings it is not available.
%
%    \begin{macrocode}
\DeclareTextCommand{\SS}{OT1}{SS}
\ProvideTextCommandDefault{\SS}{\UseTextSymbol{OT1}{\SS}}
%    \end{macrocode}
%
% \subsubsection{Shorthands for quotation marks}
%
% Shorthands are provided for a number of different quotation marks,
% which make them usable both outside and inside mathmode. 
%
% \macro{\glq}
% \macro{\grq}
%
% The `german' single quotes.
%
%    \begin{macrocode}
\ProvideTextCommandDefault{\glq}{%
  \textormath{\quotesinglbase}{\mbox{\quotesinglbase}}}
%    \end{macrocode}
%
% The definition of |\grq| depends on the fontencoding. With
% \texttt{T1} encoding no extra kerning is needed.
%
%    \begin{macrocode}
\ProvideTextCommand{\grq}{T1}{%
  \textormath{\kern\z@\textquoteleft}{\mbox{\textquoteleft}}}
\ProvideTextCommand{\grq}{TU}{%
  \textormath{\textquoteleft}{\mbox{\textquoteleft}}}
\ProvideTextCommand{\grq}{OT1}{%
  \save@sf@q{\kern-.0125em
    \textormath{\textquoteleft}{\mbox{\textquoteleft}}%
    \kern.07em\relax}}
\ProvideTextCommandDefault{\grq}{\UseTextSymbol{OT1}\grq}
%    \end{macrocode}
%
% \macro{\glqq}
% \macro{\grqq}
%
% The `german' double quotes.
%
%    \begin{macrocode}
\ProvideTextCommandDefault{\glqq}{%
  \textormath{\quotedblbase}{\mbox{\quotedblbase}}}
%    \end{macrocode}
%
% The definition of |\grqq| depends on the fontencoding. With
% \texttt{T1} encoding no extra kerning is needed.
%
%    \begin{macrocode}
\ProvideTextCommand{\grqq}{T1}{%
  \textormath{\textquotedblleft}{\mbox{\textquotedblleft}}}
\ProvideTextCommand{\grqq}{TU}{%
  \textormath{\textquotedblleft}{\mbox{\textquotedblleft}}}
\ProvideTextCommand{\grqq}{OT1}{%
  \save@sf@q{\kern-.07em
    \textormath{\textquotedblleft}{\mbox{\textquotedblleft}}%
    \kern.07em\relax}}
\ProvideTextCommandDefault{\grqq}{\UseTextSymbol{OT1}\grqq}
%    \end{macrocode}
%
% \macro{\flq}
% \macro{\frq}
%
% The `french' single guillemets.
%
%    \begin{macrocode}
\ProvideTextCommandDefault{\flq}{%
  \textormath{\guilsinglleft}{\mbox{\guilsinglleft}}}
\ProvideTextCommandDefault{\frq}{%
  \textormath{\guilsinglright}{\mbox{\guilsinglright}}}
%    \end{macrocode}
%
% \macro{\flqq}
% \macro{\frqq}
%
% The `french' double guillemets.
%
%    \begin{macrocode}
\ProvideTextCommandDefault{\flqq}{%
  \textormath{\guillemetleft}{\mbox{\guillemetleft}}}
\ProvideTextCommandDefault{\frqq}{%
  \textormath{\guillemetright}{\mbox{\guillemetright}}}
%    \end{macrocode}
%
% \subsubsection{Umlauts and tremas}
%
% The command |\"| needs to have a different effect for different
% languages. 
%
% \macro{\umlauthigh}
%
% \macro{\umlautlow}
% To be able to provide both positions of |\"| we provide two commands
% to switch the positioning, the default will be |\umlauthigh| (the
% normal positioning).
%
%    \begin{macrocode}
\def\umlauthigh{%
  \def\bbl@umlauta##1{\leavevmode\bgroup%
      \accent\csname\f@encoding dqpos\endcsname
      ##1\bbl@allowhyphens\egroup}%
  \let\bbl@umlaute\bbl@umlauta}
\def\umlautlow{%
  \def\bbl@umlauta{\protect\lower@umlaut}}
\def\umlautelow{%
  \def\bbl@umlaute{\protect\lower@umlaut}}
\umlauthigh
%    \end{macrocode}
%
% \macro{\lower@umlaut}
% Used to position the |\"| closer to the letter.
%
%    \begin{macrocode}
\expandafter\ifx\csname U@D\endcsname\relax
  \csname newdimen\endcsname\U@D
\fi
\def\lower@umlaut#1{%
  \leavevmode\bgroup
    \U@D 1ex%
    {\setbox\z@\hbox{%
      \char\csname\f@encoding dqpos\endcsname}%
      \dimen@ -.45ex\advance\dimen@\ht\z@
      \ifdim 1ex<\dimen@ \fontdimen5\font\dimen@ \fi}%
    \accent\csname\f@encoding dqpos\endcsname
    \fontdimen5\font\U@D #1%
  \egroup}
\AtBeginDocument{%
  \DeclareTextCompositeCommand{\"}{OT1}{a}{\bbl@umlauta{a}}%
  \DeclareTextCompositeCommand{\"}{OT1}{e}{\bbl@umlaute{e}}%
  \DeclareTextCompositeCommand{\"}{OT1}{i}{\bbl@umlaute{\i}}%
  \DeclareTextCompositeCommand{\"}{OT1}{\i}{\bbl@umlaute{\i}}%
  \DeclareTextCompositeCommand{\"}{OT1}{o}{\bbl@umlauta{o}}%
  \DeclareTextCompositeCommand{\"}{OT1}{u}{\bbl@umlauta{u}}%
  \DeclareTextCompositeCommand{\"}{OT1}{A}{\bbl@umlauta{A}}%
  \DeclareTextCompositeCommand{\"}{OT1}{E}{\bbl@umlaute{E}}%
  \DeclareTextCompositeCommand{\"}{OT1}{I}{\bbl@umlaute{I}}%
  \DeclareTextCompositeCommand{\"}{OT1}{O}{\bbl@umlauta{O}}%
  \DeclareTextCompositeCommand{\"}{OT1}{U}{\bbl@umlauta{U}}}
%    \end{macrocode}
%
% Finally, make sure the default hyphenrules are defined (even if
% empty).
%
%    \begin{macrocode}
\ifx\l@english\@undefined
  \chardef\l@english\z@
\fi
\ifx\l@unhyphenated\@undefined
  \newlanguage\l@unhyphenated
\fi
%    \end{macrocode}
%
% \subsection{Layout}
%
% Layout is mainly intended to set bidi documents, but there is at
% least a tool useful in general.
%
%    \begin{macrocode}
\bbl@trace{Bidi layout}
\providecommand\IfBabelLayout[3]{#3}%
%    \end{macrocode}
%
% \subsection{Load engine specific macros}
%
% Some macros are not defined in all engines, so, after loading the
% files define them if necessary to raise an error.
%
%    \begin{macrocode}
\bbl@trace{Input engine specific macros}
\ifcase\bbl@engine
  \input txtbabel.def
\or
  \input luababel.def
\or
  \input xebabel.def
\fi
\providecommand\babelfont{\bbl@error{only-lua-xe}{}{}{}}
\providecommand\babelprehyphenation{\bbl@error{only-lua}{}{}{}}
\ifx\babelposthyphenation\@undefined
  \let\babelposthyphenation\babelprehyphenation
  \let\babelpatterns\babelprehyphenation
  \let\babelcharproperty\babelprehyphenation
\fi
%</core>
%    \end{macrocode}
%
% \section{Additional support for Plain \TeX}
%
% \subsection{Not renaming \file{hyphen.tex}}
%
% As Don Knuth has declared that the filename \file{hyphen.tex} may
% only be used to designate \emph{his} version of the american
% English hyphenation patterns, a new solution has to be found in
% order to be able to load hyphenation patterns for other languages
% in a plain-based \TeX-format.
% When asked he responded:
% \begin{quote}
%   That file name is “sacred”, and if anybody changes it they will
%   cause severe upward/downward compatibility headaches.
%
%   People can have a file localhyphen.tex or whatever they like,
%   but they mustn't diddle with hyphen.tex (or plain.tex except to
%   preload additional fonts).
% \end{quote}
%
% The files \file{bplain.tex} and \file{blplain.tex} can be used as
% replacement wrappers around \file{plain.tex} and
% \file{lplain.tex} to achieve the desired effect, based on the
% \pkg{babel} package. If you load each of them with ini\TeX, you
% will get a file called either \file{bplain.fmt} or
% \file{blplain.fmt}, which you can use as replacements for
% \file{plain.fmt} and \file{lplain.fmt}.
%
% As these files are going to be read as the first thing ini\TeX\
% sees, we need to set some category codes just to be able to
% change the definition of |\input|.
%
%    \begin{macrocode}
%<*bplain|blplain>
\catcode`\{=1 % left brace is begin-group character
\catcode`\}=2 % right brace is end-group character
\catcode`\#=6 % hash mark is macro parameter character
%    \end{macrocode}
%
% If a file called \file{hyphen.cfg} can be found, we make sure that
% \emph{it} will be read instead of the file \file{hyphen.tex}. We
% do this by first saving the original meaning of |\input| (and I
% use a one letter control sequence for that so as not to waste
% multi-letter control sequence on this in the format).
%
%    \begin{macrocode}
\openin 0 hyphen.cfg
\ifeof0
\else
  \let\a\input
%    \end{macrocode}
%
% Then |\input| is defined to forget about its argument and load
% \file{hyphen.cfg} instead. Once that's done the original meaning
% of |\input| can be restored and the definition of |\a| can be
% forgotten.
%
%    \begin{macrocode}
  \def\input #1 {%
    \let\input\a
    \a hyphen.cfg
    \let\a\undefined
  }
\fi
%</bplain|blplain>
%    \end{macrocode}
%
% Now that we have made sure that \file{hyphen.cfg} will be loaded
% at the right moment it is time to load \file{plain.tex}.
%
%    \begin{macrocode}
%<bplain>\a plain.tex
%<blplain>\a lplain.tex
%    \end{macrocode}
%
% Finally we change the contents of |\fmtname| to indicate that
% this is \emph{not} the plain format, but a format based on plain
% with the \pkg{babel} package preloaded.
%
%    \begin{macrocode}
%<bplain>\def\fmtname{babel-plain}
%<blplain>\def\fmtname{babel-lplain}
%    \end{macrocode}
%
% When you are using a different format, based on plain.tex you can
% make a copy of blplain.tex, rename it and replace \file{plain.tex}
% with the name of your format file.
%
% \section{Compatibility}
%
% A couple of proxy files, for compatibility (\texttt{plain.def} and
% \texttt{switch.def}).
%
%    \begin{macrocode}
%<*plain>
\input babel.def
%</plain>
%<*kernel>
\let\bbl@onlyswitch\@empty
\input babel.def
\let\bbl@onlyswitch\@undefined
%</kernel>
%    \end{macrocode}
%
% \Finale
%
\endinput
%
% Local Variables:
% mode: doctex
% coding: utf-8
% TeX-engine: luatex
% End:
